% !TeX program = xelatex
% Compile twice from the repository root:
% xelatex -interaction=nonstopmode -halt-on-error "-jobname=gradbars-manual-v0.0.4" docs/gradbars-manual.tex
\documentclass[10pt,a4paper,fontset=fandol]{ctexart}
\usepackage[margin=18mm,top=19mm,bottom=20mm,headheight=15pt]{geometry}
\usepackage{gradbars}
\usepackage{booktabs,array,tabularx,amsmath}
\usepackage{tcolorbox}
\tcbuselibrary{listings,skins}
\usepackage{fancyhdr}
\usepackage[colorlinks=true,linkcolor=gradbarsBlue,urlcolor=gradbarsTeal,
  pdftitle={gradbars 宏包手册},pdfauthor={SuFan}]{hyperref}
\definecolor{manualCode}{HTML}{EEEEFF}
\definecolor{manualResult}{HTML}{FAFAE6}
\definecolor{manualLine}{HTML}{DCE3ED}
\definecolor{manualMuted}{HTML}{64748B}
\setlength{\parindent}{0pt}
\setlength{\parskip}{5pt}
\setcounter{tocdepth}{1}
\renewcommand{\arraystretch}{1.3}
\ctexset{
  section={format=\Large\bfseries\color{gradbarsInk},beforeskip=0pt,afterskip=10pt},
  subsection={format=\normalsize\bfseries\color{gradbarsInk},beforeskip=12pt,afterskip=4pt}
}
\pagestyle{fancy}
\fancyhf{}
\fancyhead[L]{\small\sffamily\color{manualMuted}gradbars\quad 宏包手册}
\fancyhead[R]{\small\sffamily\color{manualMuted}v0.0.4}
\fancyfoot[L]{\footnotesize\color{manualMuted}效果与源码对照 \enspace / \enspace XeLaTeX}
\fancyfoot[R]{\small\thepage}
\renewcommand{\headrulewidth}{.3pt}
\newcommand{\key}[1]{\texttt{\detokenize{#1}}}
\newcommand{\note}[1]{\begin{tcolorbox}[colback=gradbarsTrack,colframe=gradbarsTrack,
  boxrule=0pt,arc=0pt,left=3mm,right=3mm,top=2mm,bottom=2mm]
  \small #1\end{tcolorbox}}
\lstdefinestyle{manual}{
  language=[LaTeX]TeX,
  basicstyle=\ttfamily\fontsize{8}{10.6}\selectfont,
  texcsstyle=*\color{blue!70!black},
  keywordstyle=\color{green!40!black},
  commentstyle=\color{manualMuted},
  moretexcs={gradheader,gradspan,gradgroup,gradtablestyle,gradtableschema,graduseschema,gradcolumn,gradrow,gradbullet,gradintervals,graddumbbell,gradspark,gradbarscategory,gradcompare,gradrange,gradbarsloadcsv,gradbarscsvtable,gradbarscsvstyle,gradbarcsv,gradbarscsvcell,gradbarscolumn,gradbarslegend,gradbar,gradstack,gradbarsstyle,gradbarssetup,toprule,midrule,bottomrule,arraystretch},
  morekeywords={theme,max,width,height,label,precision,unit,text,blue,teal,solid,gray,percent},
  columns=fullflexible,keepspaces=true,showstringspaces=false,
  breaklines=true,breakatwhitespace=false,tabsize=2,
  aboveskip=0pt,belowskip=0pt
}
\newtcblisting{example}[2][]{
  enhanced,skin=bicolor,listing engine=listings,text side listing,
  listing options={style=manual},
  title={#2},fonttitle=\small\bfseries,
  coltitle=gradbarsInk,colbacktitle=gradbarsTrack,
  colback=manualResult,colbacklower=manualCode,colframe=manualLine,
  boxrule=.3pt,arc=0pt,lefthand width=49mm,sidebyside gap=5mm,
  left=3mm,right=3mm,top=3mm,bottom=3mm,
  before skip=8pt,after skip=8pt,
  sidebyside align=top,
  before upper={\small\setlength{\parskip}{5pt}},
  #1
}
\newtcblisting{codeonly}{
  listing only,listing engine=listings,listing options={style=manual},
  colback=manualCode,colframe=manualCode,boxrule=0pt,arc=0pt,
  left=3mm,right=3mm,top=2mm,bottom=2mm
}
\newtcblisting{coverexample}[1]{
  enhanced,skin=bicolor,listing engine=listings,text side listing,
  listing options={style=manual,basicstyle=\ttfamily\fontsize{7.5}{9}\selectfont},
  title={#1},fonttitle=\scriptsize\bfseries,
  coltitle=gradbarsInk,colbacktitle=gradbarsTrack,
  colback=manualResult,colbacklower=manualCode,colframe=manualLine,
  boxrule=.25pt,arc=0pt,lefthand width=59mm,sidebyside gap=4mm,
  left=2mm,right=2mm,top=1mm,bottom=1mm,
  before skip=3pt,after skip=3pt,sidebyside align=center,
  before upper={\small\setlength{\parskip}{2pt}}
}
\begin{document}
\thispagestyle{empty}
{\small\sffamily\color{gradbarsBlue}THE GRADBARS PACKAGE}
\vspace{7mm}

{\fontsize{46}{50}\selectfont\sffamily\bfseries\color{gradbarsInk}gradbars}\par
\vspace{4mm}
{\LARGE\bfseries\color{gradbarsInk}\LaTeX{} 数据条与可视化表格}\par
\vspace{3mm}
{\large\color{manualMuted}让数值、比较与趋势，在表格中清晰呈现。}\par
\vspace{7mm}

\begingroup
\gradbarssetup{width=30mm,precision=0}
\begin{coverexample}{渐变、圆角与目标线}
\gradbar[rounded=2pt,target=80]{68}
\end{coverexample}
\begin{coverexample}{正负双向条与误差标记}
\gradbar[min=-50,max=50,error=8,
  theme=teal]{-24}
\end{coverexample}
\begin{coverexample}{堆叠构成}
\gradstack{25,35,20}
\end{coverexample}
\begin{coverexample}{前后对比与差值}
\graddumbbell[compare label=delta]{46}{72}
\end{coverexample}
\begin{coverexample}{浮动区间与点估计}
\gradrange[range point=64]{45}{82}
\end{coverexample}
\begin{coverexample}{迷你趋势与缺失观测}
\gradspark{30,65,NA,50,82}
\end{coverexample}
\begin{coverexample}{专属表格：自动列宽、多级表头与最优值}
\begin{gradtable}[width=\linewidth,
  label width=8mm,column sep=3pt,
  min graphic width=10mm,header align=c]
  \gradcolumn{name}{方法}
  \gradcolumn[type=bar,mark=best,
    options={precision=0}]{score}{得分}
  \gradheader{\gradspan{2}{模型评估}}
  \gradrow{方法 A,82}
  \gradrow{方法 B,91}
\end{gradtable}
\end{coverexample}
\endgroup
\vspace{3mm}
{\small\color{gradbarsInk}统一接口\quad ·\quad 灵活配色\quad ·\quad 数据语义\quad ·\quad 跨页表格}\par
\vfill
{\color{manualLine}\hrule height .5pt}
\vspace{4mm}
{\small\sffamily\color{gradbarsInk}v0.0.4\quad ·\quad 用户手册与示例
  \hfill\color{manualMuted}77 个编号示例}\par
\vspace{1mm}
{\small\href{https://github.com/Tape5try/gradbars}{GitHub / Tape5try/gradbars}
  \quad ·\quad \href{https://ctan.org/pkg/gradbars}{CTAN / gradbars}}\par

\clearpage
\begingroup
\setlength{\parskip}{0pt}
\linespread{1}\selectfont
\hypersetup{linkcolor=gradbarsInk}
\makeatletter
\renewcommand*{\l@part}[2]{%
  \addpenalty\@secpenalty
  \addvspace{.7em}%
  {\bfseries\@dottedtocline{-1}{0pt}{0pt}{#1}{#2}}%
  \nobreak}
\renewcommand*{\l@section}[2]{\@dottedtocline{1}{0pt}{2em}{\bfseries #1}{#2}}
\makeatother
\fontsize{9}{11}\selectfont\tableofcontents
\endgroup
\par\vspace{4mm}
{\small\color{manualMuted}\textbf{阅读指南}\par
本手册把安装说明、参数、77 个编号示例、选型指南与常见问题集中在同一文档中。
其中包括六个应用表格，以及双栏论文、跨页长表、黑白打印和幻灯片的真实排版案例。
浅黄色区域展示实际排版效果，浅紫色区域展示生成该效果的源码。
全书按入门、图形、外观、专属表格、案例与参考六部分组织。源码中的空行用于区分示例。
示例框内的效果与源码来自同一段代码，避免两者脱节。\par
\vspace{3mm}
适用于 gradbars 0.0.4。示例数据仅用于演示，不代表真实实验结果。\par}

\clearpage
\part{快速开始与数据基础}

\section{安装与快速开始}
\subsection{安装宏包}
把 \key{gradbars.sty} 放在主文档旁边，导言区写入
\key{\usepackage{gradbars}} 即可。无需手动创建 \key{tikzpicture}，
也无需开启 shell escape。宏包依赖 TikZ、xparse、expl3、collcell（含 array）、booktabs、colortbl 和 longtable，不强制加载 ctex。
中文文档可自行使用 ctexart；本手册使用 XeLaTeX 编译。

\subsection{最小完整文档}
将下面代码保存为一个 \key{.tex} 文件，和宏包放在同一目录即可编译。
\begin{codeonly}
\documentclass{article}
\usepackage{gradbars}
\begin{document}
\gradbarssetup{max=100, width=24mm, precision=1}
\gradbar{72.4}

\gradbar{92.7}
\end{document}
\end{codeonly}

\subsection{统一设置与图形入口}
\key{\gradbarssetup{参数列表}} 设置当前作用域的默认值；
\key{\gradbar[参数列表]{数值}} 绘制一个数据条，其中可选参数只覆盖本次绘制。
数值与最大值使用十进制数字，不接受逗号、单位或计算表达式。

\begin{example}{示例 1：统一设置与单次覆盖}
\gradbarssetup{width=23mm,
  max=100,precision=1,theme=blue}
\gradbar{72.4}

\gradbar[theme=teal]{92.7}

\gradbar{81.6}
\end{example}

\subsection{作用域}
设置遵循普通 TeX 分组规则。在导言区调用会影响整篇文档；
在一对花括号中调用只影响该组。表格单元格也有自己的分组，
因此整列共用的设置应放在表格环境之前。

\begin{example}{示例 2：局部设置不会泄漏}
{\gradbarssetup{theme=gray,max=10}
 \gradbar{8}}

\gradbar{80}
\end{example}

\clearpage
\section{选型指南}
先确定读者需要比较的是数值大小、相对零的变化、组成还是不确定性，再选择条形。
同一比较组应共用范围与宽度，并在标题或表头中写明单位。

\begin{tabularx}{\linewidth}{@{}p{23mm}p{44mm}X@{}}
\toprule
\textbf{类型} & \textbf{适用情形} & \textbf{选择与说明}\\
\midrule
哑铃图 & 前后值、两组实验对比 & 用 \key{\graddumbbell} 比较位置和差距；两端必须为同一单位。\\
趋势线 & 等间隔观测的变化形状 & 用 \key{\gradspark}；跨行比较绝对波动时共用固定纵向范围。\\

普通条 & 准确率、数量、耗时、完成量 &
用 \key{\gradbar} 表示一个非负值。越长表示越大；耗时越低越好，不能把条长直接理解为优劣。\\
正负条 & 增减幅度、差值、偏差 &
用 \key{min<0} 和正的 \key{max} 显示零点两侧；同列必须共用同一个零点位置。\\
堆叠条 & 互斥类别的构成、资源分配 &
用 \key{\gradstack}；同单位、可相加的数据可分别向零点两侧累计。不同单位或重复计数不应拼成同一组成条。\\
误差线 & 点估计及不确定区间 &
在普通条或正负条上设置 \key{error} 或 \key{error minus}/\key{error plus}。
在表注中明确区间是标准差、标准误还是置信区间。\\
\bottomrule
\end{tabularx}

\subsection{组合使用时的语义}
\key{target} 是一个参考值，\key{band} 是预先给定的参考范围，
它们本身不表示估计误差。误差参数表示从点估计向两侧延伸的非负距离，
不是区间的绝对端点；例如估计 60、区间 $[55,68]$，应设置
\key{error minus=5,error plus=8}。

若希望比较不同方案的总量，保留共同的 \key{max}；若要比较百分比组成，
先由数据处理步骤计算各段百分比，再用统一的 \key{max=100}。
仅设置 \key{segment labels=percent} 会改变标签，不会把原始数据条变成满条。

\subsection{窄列与打印的取舍}
双栏论文优先使用简短数值标签和共享图例；长类别名可放在左侧文本列。
黑白打印同时保留名称、数值或目标位置，避免仅靠颜色表达类别。
幻灯片应增大轨道高度与字体；标签外移后，应为其留出额外垂直空间。

\clearpage
\section{数值语义与边界}
默认范围为非负数据，条形从零开始，长度遵循
\[
  L=W\min\!\left(1,\frac{v}{M}\right),
  \qquad v\geq 0,\quad M>0,
\]
其中 $v$ 是输入值，$M$ 是 \key{max}，$W$ 是 \key{width}。
标签的舍入与格式不改变条长。

\begin{example}{示例 3：原始数值、单位与计算占比}
\gradbar[max=200]{50}

\gradbar[max=200,
  unit={\,ms}]{50}

\gradbar[max=200,
  value format=percent]{50}
\end{example}

\key{unit} 只添加后缀，不换算数值。
\key{value format=percent} 才计算 $100v/M$，自动添加百分号，并忽略 \key{unit}。
如果原始数据已经是 86.5 这样的百分数，应使用 \key{max=100,unit={\%}}。

\begin{example}{示例 4：已有百分数与自定义标签}
\gradbar[unit={\%}]{86.5}

\gradbar[max=10,
  text={8 / 10}]{8}

\gradbar[max=10,
  text={},precision=0]{8}
\end{example}

\begin{example}{示例 5：零值和超出最大值}
\gradbar{0}

\gradbar{100}

\gradbar[max=100,
  value format=percent]{125}
\end{example}

零值只有空轨道，不会被强制画出最短长度。很小的正值也不会人为放大；
实际可见精度受 TeX 长度分辨率限制。
超出最大值时默认发出警告，并截断彩色区域，标签仍显示真实值。

\note{\key{overflow=error} 会把越界改成明确报错。
默认范围不接受负数；需要双向条时设置负的 min。非数字、零或负的最大值同样会报错。}

\clearpage
\section{命名样式与缺失值}
\key{\gradbarsstyle{名称}{参数}} 定义可复用的配置；使用 \key{style=名称} 调用。
样式可以引用已定义的样式，后面的参数覆盖前面的设置。
定义和调用均遵循普通 TeX 分组，不会改变其他单元格。

\begin{example}{示例 6：定义并复用样式}
\gradbarsstyle{accuracy}{
  max=100,theme=teal,
  rounded=2pt,precision=1,
  unit={\%}
}
\gradbar[style=accuracy]{92.7}

\gradbar[style=accuracy,
  target=95]{94.2}
\end{example}

样式在当前作用域内可以重新定义。未知名称会报错；循环引用会触发嵌套深度限制。

\begin{example}{示例 7：零值与缺失值}
\gradbar{0}

\gradbar{NA}

\gradbar[missing text={N/A}]{}
\end{example}

空参数或大写 \key{NA} 表示缺失值，忽略首尾空格。
零值显示数值零；缺失值只显示轨道和缺失标签，默认标签为短横线。
宏包不会把其他文字或拼写错误自动当成缺失值。

\note{缺失值不绘制数值填充和误差线，也不附加单位或百分号。
目标线和区间背景仍保留，表示这一列的参考尺度。
\key{text} 可以覆盖缺失标签，\key{label=none} 可以隐藏它。}

\clearpage
\section{数字格式与区间背景}
\key{value format} 决定显示原值还是计算占比；
\key{number format} 决定数字的排版形式，不改变条长。
\key{fixed} 为默认定点格式，千位使用细空格；\key{grouped} 使用逗号分组；
\key{scientific} 使用科学计数法。\key{precision} 指定小数位数，科学计数法中指尾数的小数位数。

\begin{example}{示例 8：分组数字与科学计数法}
\gradbar[max=200000,
  number format=grouped,
  precision=0]{125000}

\gradbar[max=1,
  number format=scientific,
  precision=2]{0.00001234}
\end{example}

科学计数法只用于输出，输入仍使用十进制数。
极小数不会先按定点精度四舍五入成零；标签较长时应增大 \key{label width}。

\begin{example}{示例 9：参考区间与目标线}
\gradbar[band={60,80},
  band color=lightgreen,
  target=70]{75}

\gradbar[min=-50,max=50,
  band={-10,15},
  band color=skyblue,
  band opacity=0.4]{-25}
\end{example}

\key{band={下限,上限}} 使用与数据相同的单位，要求两个端点严格递增且位于量程内。
背景在轨道上下各延伸 2pt，填充经过区间时仍可看见参考范围。
\key{band color} 默认为 gold，\key{band opacity} 默认为 0.3，取值为 0 到 1。
使用 \key{band={}} 清除区间。

\note{区间背景表示人为给定的参考范围，不代表统计误差。
误差区间仍使用 \key{error}；参考区间的越界端点会直接报错，不自动截断。}

\clearpage
\part{图形与比较方式}

\section{正负双向条与误差标记}
\key{min} 默认是零。设置负的下限后，零线位于
$W(-\mathrm{min})/(\mathrm{max}-\mathrm{min})$，
负值向左、正值向右。非对称范围的零线不在轨道中心。
上限必须为正，下限必须不大于零。

\begin{example}{示例 10：共用零线的双向条}
\gradbarssetup{
  min=-50,max=50,
  width=30mm,
  negative color=lightred
}
\gradbar{-30}

\gradbar{0}

\gradbar{30}
\end{example}

负值默认使用 \key{negative color=lightred} 的纯色；
正值使用所选主题。跨零范围不接受 \key{value format=percent}，
避免把位置比例当作有意义的百分比；百分数变化可直接使用 \key{unit}。

\begin{example}{示例 11：对称与非对称误差}
\gradbar[error=5]{72}

\gradbar[error minus=4,
  error plus=9]{72}

\gradbar[min=-50,max=50,
  error minus=8,error plus=12]{-5}
\end{example}

\key{error=e} 表示区间 $[v-e,v+e]$；
\key{error minus} 和 \key{error plus} 分别指定非负的上下误差幅度，
不是区间端点。只指定一侧时，另一侧默认为零。
误差单位与数据相同，\key{error color} 控制线色；
\key{error={}} 清除两侧误差。误差含义由用户说明，宏包不计算置信区间。

\begin{example}{示例 12：组合效果}
\gradbar[
  min=-50,max=50,width=30mm,
  rounded=2pt,target=20,
  error=5,label=end
]{-25}
\end{example}

\note{数值及误差端点超出范围时，默认发出警告并截断显示，
原始标签保持不变。截断的误差端点不代表真实端点。
严谨的区间图建议使用 \key{overflow=error}，或扩大范围以显示完整误差。}

\clearpage
\section{堆叠条}
\key{\gradstack[参数]{数值列表}} 将正数向右累计、负数向左累计。含负段时必须设置负的 min。
所有段共用 \key{max}，默认 100；不会自动把总量归一化成满条。
默认总标签显示代数和（净值），也可以用 \key{value format=percent} 显示总量占上限的比例。

\begin{example}{示例 13：三段组成与未填满的轨道}
\gradstack[
  stack colors={skyblue,rose,gold},
  rounded=2pt,
  target=90
]{30,25,15}

\gradstack[max=200,
  value format=percent,
  stack colors={skyblue,rose,gold}
]{30,25,15}
\end{example}

第一条的三段长度分别占轨道的 30\%、25\%、15\%，总量为 70；
第二条沿用相同原始数据，上限改为 200，总量占比为 35\%。
\key{stack colors} 依次指定段颜色；颜色少于段数时循环使用。
零段不占宽度，但仍保留其颜色序号。所有段均为零时只显示空轨道。

\begin{example}[text and listing,sidebyside=false]{示例 14：表格中的堆叠组成}
\gradbarsstyle{composition}{
  width=42mm,precision=0,rounded=2pt,
  stack colors={skyblue,rose,gold}
}
\begin{tabular}{@{}lccc l@{}}
  \toprule
  Plan & A & B & C & Total / 100 \\
  \midrule
  Alpha & 30 & 25 & 15 & \gradstack[style=composition]{30,25,15} \\
  Beta  & 20 & 35 & 35 & \gradstack[style=composition]{20,35,35} \\
  \bottomrule
\end{tabular}
\end{example}

\note{堆叠条不接受缺失段、空列表、条件配色或误差线。含负段时显式设置 \key{min<0}。
目标线、区间背景、圆角、数字格式和标签位置可以组合使用。
正负小计分别检查范围，越界时警告并截断相应端，标签保留真实数据；\key{overflow=error} 改为报错。}


\subsection{正负混合堆叠与小计}
正数和负数分别从零开始累计，输入顺序只决定同一侧的段顺序。
净值为零不代表没有活动量；例如 50 与 $-50$ 仍绘制两段。
默认总标签为净值；\key{stack totals=separate} 显示“正小计 / 负小计”。
\begin{example}[text and listing,sidebyside=false]{示例 15：收入与支出分开累计}
\gradstack[min=-60,max=100,width=85mm,
  height=12pt,precision=0,palette=categorical,
  stack totals=separate,
  stack names={Sales,Materials,Service,Operations},
  segment labels=value,legend=true]{60,-25,20,-15}
\end{example}
分段百分比使用 $100|v_i|/\sum_j|v_j|$，表示绝对活动量份额，所有非零段份额相加为 100\%。
负号由数值标签或左右位置表达。全零数据没有可见段，也不计算百分比。
这不同于净值占比：净值可为零，不能拿它作为组成比例的分母。
负量程不支持总标签的 \key{value format=percent}；需要百分号单位可使用 \key{unit={\%}}。

\clearpage
\section{堆叠图例与分段标签}
\key{stack names} 为每一段提供名称，顺序与数值、颜色一致。
\key{legend=true} 在条形下方自动生成图例；名称较长时在轨道宽度内换行。
\key{segment labels} 可显示数值、占比、名称或它们的组合。

\begin{example}[text and listing,sidebyside=false]{示例 16：自动图例与小段标签外移}
\gradstack[
  width=70mm,height=16pt,precision=0,
  stack colors={gradbarsBlue,gradbarsTeal,gold},
  stack names={Train,Validation,Test},
  segment labels=name percent,legend=true,
  label=end
]{80,2,18}
\end{example}

段内放不下的标签自动移至条形上方，并添加引导线；水平范围相交的外部标签逐层抬高。\par
\key{segment text color=auto} 默认按段颜色选择黑字或白字。
\key{segment font} 默认是 \key{\scriptsize}，可按版面修改。
零值段不绘制标签，但仍保留图例项和配色序号。

\begin{example}[text and listing,sidebyside=false]{示例 17：分段占比与总量占比采用不同分母}
\gradstack[
  max=200,width=80mm,height=16pt,precision=0,
  stack colors={gradbarsBlue,gradbarsTeal},
  segment labels=percent,value format=percent
]{20,30}
\end{example}

这里总量为 50，条长占量程 200 的四分之一，总标签为 25\%。
两段标签分别是 $20/50=40\%$ 和 $30/50=60\%$。
分段占比以原始段值的绝对值之和为分母；此处均为正数，因而等于原始总量。
即使条长被截断，也不会按可见部分重新归一化。
\key{segment labels=value} 显示原始分段值，沿用数字格式及精度，省略总标签的 \key{unit}。

\begin{example}[text and listing,sidebyside=false]{示例 18：多个堆叠条共用一个图例}
\gradbarsstyle{parts}{width=55mm,height=14pt,
  precision=0,stack colors={skyblue,rose,gold}}
\begin{tabular}{@{}ll@{}}
  A & \gradstack[style=parts]{50,30,20}\\[3pt]
  B & \gradstack[style=parts]{65,20,15}
\end{tabular}

\gradbarslegend[style=parts]{Train,Validation,Test}
\end{example}

\note{启用自动图例或包含名称的段标签时，必须提供与段数一致的 \key{stack names}。
\key{segment labels} 可选 none、value、percent、name、name value、name percent。
独立图例需要和数据条使用相同的 \key{stack colors}，建议通过命名样式统一设置。}

\clearpage
\section{双层对比条、棒棒糖与浮动区间}
这三种图形沿用现有的 min、max、width、标签格式与命名样式。
比较时应保持相同的范围与宽度；图形形状不改变原始数据的意义。

\subsection{双层对比条：基准与当前值}
\key{\gradcompare[选项]{基准值}{当前值}} 绘制宽而浅的基准条，
并在其中央叠加较细的当前值条。两个值分别从零开始定位，不相加，也不做除法。
默认标签按“当前值 / 基准值”排列；斜线是分隔符，不表示比率。

\begin{example}[text and listing,sidebyside=false]{示例 19：共用尺度的双层比较}
\gradcompare[width=65mm,height=12pt,
  compare color=black!20,precision=0]{60}{82}

\gradcompare[width=65mm,height=12pt,
  compare ratio=.35,precision=0]{85}{65}

\gradcompare[min=-50,max=100,width=65mm,
  height=12pt,precision=0]{-20}{35}
\end{example}

\key{compare color} 默认是 \key{black!20}，仅设置基准层颜色。
\key{compare ratio} 是当前层高度与总高度的比值，默认 0.45，必须严格位于 0 与 1 之间。
当前层沿用主题、条件配色和纹理；目标线与误差线也仍针对当前值。
负值要求显式设置负的 min。任一值越界时，沿用 clip 警告并截断或 error 报错规则，
标签保留原始数值。

基准与当前值可分别使用空值或 NA 表示缺失，对应层不绘制，标签显示缺失标记。
总标签放在外部或条尾上方；请求 inside/auto 时会移至外部，避免遮挡两层之间的差异。
需要简短标签时，可用 text 覆盖，但应在表头或表注中说明基准值。

\clearpage
\subsection{棒棒糖：数值位置与细线}
\key{shape=lollipop} 直接用于 \key{\gradbar}，无需换用另一套数据接口。
细线从零延伸至数值位置，圆点标出数值；浅色细线表示整个轨道。
适合密集表格中需要减少填色面积的场景。

\begin{example}[text and listing,sidebyside=false]{示例 20：棒棒糖的正负值、零值与缺失值}
\gradbarssetup{shape=lollipop,min=-50,max=100,
  width=65mm,marker size=2pt,stem width=.7pt,
  precision=0}
\gradbar{72}

\gradbar{-25}

\gradbar{0}

\gradbar{NA}
\end{example}

\key{marker size} 指圆点半径，默认 2pt，并限制在轨道高度的一半以内。
\key{stem width} 默认 0.6pt，控制棒棒糖细线及浮动区间端点线的粗细。
零值在零点绘制圆点，缺失值不绘制圆点或数值细线，两者保持区别。
端点圆点可以向轨道外伸出一个半径，应给表格边界留出空间。
条件配色仍由真实数值决定；使用纯色，不使用渐变或面积纹理。
inside/auto 标签移至外部，end 标签仍可放在端点上方。

\subsection{浮动区间：直接填写上下界}
\key{\gradrange[选项]{下界}{上界}} 在两个端点之间绘制区间条，
默认标签显示原始端点。\key{range point} 可显式提供一个区间内的点，以空心圆标出；
宏包不会自行计算均值、中位数或置信区间。

\begin{example}[text and listing,sidebyside=false]{示例 21：浮动区间与区间内的点}
\gradrange[width=65mm,height=9pt,
  range point=50,precision=0]{35}{75}

\gradrange[min=-50,max=100,width=65mm,
  height=9pt,precision=0]{-20}{35}

\gradrange[width=65mm,height=9pt,precision=0]{50}{50}

\gradrange[width=65mm,height=9pt]{NA}{NA}
\end{example}

\clearpage
\subsection{区间语义与组合规则}
下界必须小于或等于上界；相等时绘制端点线，不制造虚假的区间宽度。
上下界必须同时有值或同时缺失。range point 必须位于原始上下界之间，缺失区间不能包含点。
上下界越界时沿用 overflow 规则；所有位置使用同一 min/max 映射。
范围下限仍要求为零或负值，上限仍要求为正值。

浮动区间本身就是数据区间，不接受 error 或 thresholds，以免把上界误当作点估计或评价对象。
可以添加 target、band、圆角和纹理。
添加 range point 后，内部总标签自动外移，避免盖住圆点；默认仍显示上下界，而不是点值。
value format=percent 分别计算两个端点相对于 max 的百分比，仅用于 min=0。

\begin{tabularx}{\linewidth}{@{}>{\raggedright\arraybackslash}p{30mm}X@{}}
\toprule
\textbf{需要表达的信息} & \textbf{推荐方式}\\
\midrule
当前值与基准值 & gradcompare：两条共用零点与尺度，不累加。\\
单个数值的大小 & gradbar 或 shape=lollipop：填充长度或端点位置表示数值。\\
已知的两个区间端点 & gradrange：直接填写上下界，可附加区间内的点。\\
点估计与误差距离 & gradbar 的 error/error minus/error plus：填写相对点估计的距离。\\
所有观测共用的参考区间 & band：绘制背景参考，不充当单条观测的区间。\\
\bottomrule
\end{tabularx}

\clearpage
\section{哑铃对比图}
\key{\graddumbbell[选项]{基准值}{当前值}} 用空心圆表示基准、实心圆表示当前值，连接线只连接两个观测位置。
它与双层条共用范围、数值格式、目标与误差参数。误差线针对当前值。
默认标签按“当前值 / 基准值”排列；\key{compare label=delta} 显示当前值减去基准值。
\begin{example}[text and listing,sidebyside=false]{示例 22：位置比较与差值}
\graddumbbell[width=65mm,palette=categorical,
  marker size=2.5pt,precision=0]{48}{76}

\graddumbbell[width=65mm,min=-50,max=100,
  compare label=delta,precision=0]{-20}{35}

\graddumbbell[width=65mm,precision=0]{NA}{65}
\end{example}
缺失一端只绘制另一端，连接线不绘制；差值标签在任一端缺失时显示缺失标记。
两端相同会重合，默认数值标签仍保留两个原始值。内部标签自动外移。
使用 \key{compare color} 设置基准圆颜色，\key{color} 设置当前圆及连接线颜色。
\key{marker size} 和 \key{stem width} 控制圆点半径与线宽。
\key{compare label=delta} 同样可用于双层条。
哑铃与趋势线的专用命令会保留各自图形类型，命名样式中的排版预设不会将它们改回普通条。

\clearpage
\section{对比变化标注}
\key{compare annotation=none|arrow|bracket} 在双层条或哑铃上方添加方向箭头或差值括号。
默认 none。箭头从基准指向当前值；相等时省略箭头，仍可显示零差值。
缺失任意一个值时不画变化标注。

\begin{example}[text and listing,sidebyside=false]{示例 23：差值括号与百分比变化}
\gradcompare[width=65mm,precision=0,
  compare annotation=bracket]{80}{55}

\graddumbbell[width=65mm,precision=1,
  compare annotation=arrow,compare label=percent,
  marker shape=diamond]{40}{60}

\graddumbbell[width=65mm,precision=1,
  compare label=percent]{0}{30}
\end{example}
\key{compare label=percent} 计算
\[
  100\,\frac{\text{当前值}-\text{基准值}}{|\text{基准值}|}\,\% .
\]
这是相对基准的变化，区别于 value format=percent 的量程占比。
基准为零或任一值缺失时，百分比标签显示缺失符号，不生成无穷大或虚构的百分比。
负基准使用绝对值作分母；正差值加正号，负差值保留负号。
compare annotation 的文本跟随 percent 模式，其余模式显示原单位差值。
标注与参考名称可能占用同一上方区域，密集情况建议保留其中一组名称，或增大行距。

\clearpage
\section{迷你趋势线}
\key{\gradspark[选项]{观测列表}} 按输入顺序绘制等间隔数据，默认标签显示最后一个位置的原始值。
空字段或 NA 保留横向位置并断开连接；末项缺失时标签也显示缺失，不回退到之前的值。
\begin{example}[text and listing,sidebyside=false]{示例 24：缺失断点与共享纵向范围}
\gradspark[width=65mm,height=6mm,precision=0,
  spark range=fixed,min=0,max=100]{25,45,30,NA,65,80}

\gradspark[width=65mm,height=6mm,precision=0,
  spark range=fixed,min=0,max=100]{45,42,50,60,65,70}

\gradspark[width=65mm,height=6mm,precision=0]{5,5,5}
\end{example}
\key{spark range=auto} 默认使用本行最小值与最大值，只适合看形状；不同自动范围的行不能比较绝对波动。
恒定序列居中，单点居中绘制；全缺失序列仅显示轨道与缺失标签。
\key{spark range=fixed} 使用统一的 min/max，遵循越界警告截断或报错规则；固定范围仍要求 min 不大于零、max 为正。
默认高度为 4ex，可覆盖。自动范围不支持计算百分比标签。

\key{spark points=extrema} 默认标出全部并列极值及末点；另可选择 all、last、none。
最低点用 \key{spark low color}，最高点用 \key{spark high color}，两者相同时使用最高点颜色。
末点若不是极值则使用线条颜色。点色表示位置特征，不表示指标优劣。
\key{color}、\key{stem width}、\key{marker size} 设置线色、线宽与点半径。
趋势线不使用面积纹理、柱形目标线、误差线、区间背景或阈值评价；标签始终在外部。

\clearpage
\section{多重参考与区间对比}
本节介绍多条参考线、子弹图和多组区间对比。所有示例数据仅用于演示；参考标记和背景不改变原始数值及条形尺度。

\subsection{多条参考线}
\key{references} 接受一组记录，每条记录按
\key{数值/颜色/线型/名称} 填写，用逗号分隔并分别加花括号。
线型可选 \key{solid}、\key{dashed}、\key{dotted}。
数值必须在 min/max 内；名称可留空，但最后一个斜杠必须保留。
名称交替采用两层高度，长名称或密集参考仍需增加 width，或在表头解释名称。

\begin{example}[sidebyside=false]{示例 25：基准、目标和警戒值}
\gradbar[width=85mm,height=9pt,precision=0,
  references={{40/gray/dotted/基准},
              {75/gradbarsTeal/dashed/目标},
              {95/gradbarsOrange/solid/警戒}}]{68}
\end{example}
\key{references={}} 清除继承的参考线；原有 \key{target} 仍可共存。
参考线适用于普通条、正负条、堆叠、双层条、哑铃及区间条，不用于趋势线。
参考名称不参与数值格式化；名称内不要使用作为字段分隔符的斜杠。

\subsection{子弹图：分级背景、实际值和目标}
\key{\gradbullet[选项]{实际值}} 绘制较高的背景和居中的细条。
必须提供 \key{bullet bands}，每条记录为 \key{上界/颜色}。
第一个区间从 min 开始，后续区间从前一个上界开始；上界必须严格递增，最后一个必须等于 max。
背景区间只是参考等级，不会自动改变实际值颜色，也不会自动判定优劣方向。

\begin{example}[sidebyside=false]{示例 26：完成度子弹图}
\gradbullet[width=90mm,max=120,precision=0,
  color=gradbarsBlue,
  bullet bands={{60/black!8},{90/black!16},{120/black!26}},
  references={{90/gradbarsTeal/dashed/目标},
              {110/gradbarsOrange/solid/挑战}}]{82}
\end{example}
默认总高为 12pt，实际条占高度的 40\%，上下居中。
外部数值保留原值；NA 表示缺失，零表示实际值为零。
可使用原有条件配色、目标线和误差标记。
背景采用显式颜色；黑白打印时请像上例一样选用不同灰度。
\key{bullet bands} 仅用于子弹图，背景等级建议在表头或正文中说明。

\clearpage
\subsection{多组区间对比}
\key{\gradintervals[选项]{记录列表}} 将多组区间排列在同一个单元格内。
每条记录按 \key{名称/下界/点估计/上界/颜色} 填写。
所有行共用 width、min、max；左侧名称自动按最长名称对齐。
点估计为 NA 时不画圆点，区间仍保留；上下界同时为 NA 表示整组缺失。

\begin{example}[sidebyside=false]{示例 27：同一指标的三组不确定区间}
\gradintervals[width=72mm,min=0,max=100,
  height=5pt,precision=0,row padding=4pt,
  references={{70/gray/dashed/}}]{
  {方法 A/45/58/72/gradbarsBlue},
  {方法 B/55/68/82/gradbarsTeal},
  {方法 C/62/NA/88/gradbarsOrange}}
\end{example}
圆点为显式提供的点估计，外部标签为原始上下界；不从上下界推断均值或置信水平。
若区间表示 95\% 置信区间，应在表头或图注中写明。
点估计必须位于原始区间内，下界不得大于上界。
越界处理继承 \key{overflow}，所有组必须采用相同单位。

\begin{example}[sidebyside=false]{示例 28：正负区间与缺失组}
\begin{tabular}{ll}
指标 & 估计区间\\
变化量 & \gradintervals[min=-40,max=60,width=65mm,
  height=4pt,precision=0,row padding=3pt]{
  {组 A/-25/-10/5/gradbarsBlue},
  {组 B/5/18/35/gradbarsTeal},
  {组 C/NA/NA/NA/gray}}\\
\end{tabular}
\end{example}
该命令使用内部小表格，整体作为一个单元格排版，不在组间自动分页。
名称列、图区及外部区间标签都占用宽度；长名称请缩写，或调小图区 width。
参考名称会在每组重复，紧凑表格推荐留空，并在共同表头解释参考线。
记录名称及颜色字段不应包含斜杠；带逗号的名称应加额外花括号保护。

\clearpage
\part{外观、标签与评价规则}

\section{主题与内置颜色}
\key{theme} 一次设置填充方式、渐变两端、轨道和文字颜色。
主题参数之后仍可逐项覆盖颜色；同一选项列表从左向右执行。
示例 29 集中展示 16 种内置配色，包括四种基础主题、三种渐变预设和九种单色。
渐变预设通过 \key{theme} 调用，内置颜色通过 \key{color} 设置纯色填充。

\begin{example}[before upper={\fontsize{9}{10.6}\selectfont\setlength{\parskip}{13.8pt}}]{示例 29：内置主题与颜色}
\gradbar[theme=blue]{72}

\gradbar[theme=teal]{72}

\gradbar[theme=solid]{72}

\gradbar[theme=gray]{72}

\gradbar[theme=lbyellow]{72}

\gradbar[theme=viblue]{72}

\gradbar[theme=cyblu]{72}

\gradbar[color=lightgreen]{72}

\gradbar[color=lightyellow]{72}

\gradbar[color=lightblue]{72}

\gradbar[color=lightred]{72}

\gradbar[color=rose]{72}

\gradbar[color=skyblue]{72}

\gradbar[color=gold]{72}

\gradbar[color=lavender]{72}

\gradbar[color=peach]{72}
\end{example}

\clearpage
\begingroup
\small
\renewcommand{\arraystretch}{1.15}
\begin{tabularx}{\linewidth}{@{}lX@{}}
\toprule
\textbf{主题或颜色} & \textbf{特征与适用场景}\\
\midrule
\key{blue} & 默认蓝色渐变，适合通用表格。\\
\key{teal} & 青绿色渐变，适合区分另一组指标。\\
\key{solid} & 蓝色纯色填充，减少渐变带来的视觉干扰。\\
\key{gray} & 灰度填充，适合黑白输出；仍应保留数值标签。\\
\midrule
\key{lbyellow} & 浅蓝至浅黄渐变，适合浅色背景下的柔和展示。\\
\key{viblue} & 浅紫至浅蓝渐变，适合低饱和度的对比表格。\\
\key{cyblu} & 青色至浅蓝渐变，适合突出一组重点指标。\\
\midrule
\key{lightgreen} & 浅绿纯色，适合展示完成度或达标指标。\\
\key{lightyellow} & 浅黄纯色，适合轻量提示；浅色背景下应保留深色标签。\\
\key{lightblue} & 浅蓝纯色，适合通用数据和大面积表格。\\
\key{lightred} & 浅红纯色，适合标出需要关注的指标。\\
\key{rose} & 玫红纯色，适合强调重点数据或区分对照组。\\
\key{skyblue} & 天蓝纯色，适合清晰呈现进度和比例。\\
\key{gold} & 金黄纯色，适合突出关键结果或优选方案。\\
\key{lavender} & 淡紫纯色，适合柔和区分不同类别。\\
\key{peach} & 蜜桃纯色，适合暖色风格的辅助指标。\\
\bottomrule
\end{tabularx}
\endgroup

\begin{example}{示例 30：自定义渐变与轨道}
\gradbar[
  theme=blue,
  left color=violet!80!blue,
  right color=violet!20,
  track color=violet!5
]{68}

\gradbar[
  theme=teal,
  fill=solid,
  left color=orange!85!black,
  track color=orange!10
]{68}
\end{example}

\subsection{颜色参数}
内置配色无需自行定义。主题和渐变预设使用 \key{theme}；单色使用 \key{color}。
单色名称也可用于 \key{left color} 或 \key{right color}，名称不带反斜杠。

\key{theme=blue,color=rose} 使用玫红纯色；
\key{color=rose,theme=blue} 则恢复蓝色渐变。

\key{left color} 是渐变起点，也决定纯色填充的颜色；
\key{right color} 仅在 \key{fill=gradient} 时生效。
\key{track color} 控制空轨道，\key{text color} 控制标签。
颜色可使用 xcolor 的混色表达式。

\note{不要只靠颜色传递关键含义。并列指标应在列标题中写明名称和单位，
并保留可读的数值。可设置 \key{text color=auto} 根据背景选择黑白文字，也可显式指定颜色。}

\clearpage
\section{标签、尺寸与正文内嵌}
\begin{example}{示例 31：标签位置}
\gradbar[label=outside]{64}

\gradbar[label=inside,
  height=2.5ex,
  right color=white]{64}

\gradbar[label=none]{64}
\end{example}

外部标签采用右对齐的槽位，默认在文字过长时扩展；同列应统一设置足够大的槽宽。
内部标签位于整条轨道右端附近，不跟随彩色条的末端移动。
\key{label=none} 则只显示轨道和填充。

\begin{example}{示例 32：单位、精度和字体}
\gradbar[max=5000,
  width=18mm,label width=6em,
  precision=0,unit={\,ms}]{1250}

\gradbar[width=18mm,
  precision=2,
  font={\small\bfseries}]{86.537}
\end{example}

\begin{example}{示例 33：宽度与高度}
\gradbar[width=18mm,
  height=1ex]{70}

\gradbar[width=26mm,
  height=2ex]{70}

\gradbar[width=26mm,
  label gap=1em,
  label width=3em]{70}
\end{example}

\key{width} 只指定轨道宽度。带外部标签时，总宽度约为轨道宽度、
\key{label gap} 和 \key{label width} 三者之和。
这些长度支持 mm、pt、em、ex 等单位。

\begin{example}{示例 34：直接放入正文}
Completion:
\gradbar[width=15mm,
  height=1.2ex,
  label width=3em,
  unit={\%},precision=0]{64}

The task is progressing.
\end{example}

\note{标签默认沿用周围文字的字体。普通表格行高仍由表格本身控制；
需要更宽松的行距时设置 \key{\arraystretch}。宏包不会修改整张表格的行距。}

\clearpage
\section{智能标签与数值列}
\subsection{自动黑白文字、长标签与碰撞避让}
\key{label=auto} 优先将总标签放在已填充部分的中央，空间不足时移至轨道右侧。
\key{text color=auto} 根据标签中心处的背景选择黑白文字；渐变按中心颜色估算。
主题会设置文字颜色，因此把 \key{text color=auto} 放在 \key{theme} 后面。

\begin{example}[text and listing,sidebyside=false]{示例 35：深浅背景与长标签}
\gradbar[color=black,width=45mm,height=15pt,
  label=auto,text color=auto,precision=0]{90}

\gradbar[color=gold,width=45mm,height=15pt,
  label=auto,text color=auto,precision=0]{90}

\gradbar[width=25mm,height=15pt,label=inside,
  text={A label wider than the track}]{60}
\end{example}

默认 \key{label overflow=auto} 会将过宽的内部或条尾标签外移，
并按实际文字宽度扩展外部标签槽；过高的内部标签也会外移。
\key{label overflow=allow} 允许手工布局时保留越界标签。
目标线穿过内部标签时标签外移；带误差线的内部总标签移至右侧。
堆叠条启用段标签后，内部总标签也移至右侧，条尾总标签则位于外移段标签上方。

\note{避让发生在同一数据条内部，不会重新排版相邻单元格或相邻数据条。
外部文字默认按浅色页面选择黑字；深色页面请显式指定文字颜色。
自动颜色基于标签中心，不能保证跨越多种颜色的长标签处处对比度相同。}

\clearpage
\subsection{单元格只填写数字}
\key{\gradbarscolumn{G}{...}} 定义一个收集单元格内容的数值列。
选择尚未使用的单个英文字母作为列名，整列共用范围、样式与标签格式。
普通文本表头必须用 \key{\multicolumn} 绕过数值解析。

\begin{example}[text and listing,sidebyside=false]{示例 36：统一列样式、缺失值和零值}
\gradbarsstyle{score}{max=100,width=35mm,precision=0}
\gradbarscolumn{G}{style=score}
\begin{tabular}{@{}lG@{}}
  \toprule
  Model & \multicolumn{1}{c}{Score / 100}\\
  \midrule
  Alpha & 82\\
  Beta & 65\\
  Pending & NA\\
  Empty & \\
  Zero & 0\\
  \bottomrule
\end{tabular}
\end{example}

该接口基于 array 与 collcell，适用于这里演示的 tabular 和 longtable。
表格命令、文本、单位或已绘制的数据条不能直接作为 G 列的数值输入；
需要特殊单元格时使用 \key{\multicolumn{1}{l}{...}}。
它不会扫描整列或自动计算最大值，也不是 siunitx 的 S 列。

\clearpage
\section{圆角、目标线与条尾标签}
\begin{example}{示例 37：圆角与真实长度}
\gradbar[rounded=2pt]{72}

\gradbar[rounded=10pt]{1}

\gradbar[rounded=10pt]{0}
\end{example}

\key{rounded} 为圆角半径，默认 \key{0pt}。实际半径会被限制在轨道高度和
彩色条宽度允许的范围内。小数值不被放大，零值仍只显示空轨道。

\begin{example}{示例 38：目标线}
\gradbar[target=80]{72}

\gradbar[target=80,
  target color=red!70!black]{92}
\end{example}

\key{target} 使用与输入值相同的单位，不是相对位置。目标线必须位于
\key{min} 到 \key{max} 之间，越界会报错。
默认不绘制目标线；\key{target={}} 可清除局部或继承的目标设置。

\begin{example}{示例 39：标签跟随彩色条末端}
\gradbar[label=end]{15}

\gradbar[label=end]{72}

\gradbar[label=end]{100}
\end{example}

\key{label=end} 把标签放在彩色条末端上方。靠近左右边界时，
会在轨道范围内向内调整标签位置，避免常规数值标签越界。
因此它会增加行高，但不为外部标签槽预留水平空间。

\note{很长的自定义标签可能比轨道本身还宽；此时应增大 \key{width}，
或改用 \key{label=outside} 并设置足够的 \key{label width}。
上方标签不改变正文中的轨道基线。}

\clearpage
\section{越界指示与点标记}
\key{overflow mark=arrow|slash|none} 控制量程边缘的箭头、双斜线或不标记。
单图默认 none，gradtable 默认 arrow。只有严格超出端点时显示；刚好等于上限不显示。
\key{overflow=error} 仍会报错，标记不能绕过校验；裁剪模式保留真实标签。

\begin{example}[text and listing,sidebyside=false]{示例 40：超范围数值和区间}
\gradbar[max=100,width=65mm,precision=0,
  overflow mark=arrow]{125}

\gradbar[min=-40,max=60,width=65mm,precision=0,
  overflow mark=slash]{-55}

\gradrange[min=-40,max=60,width=65mm,precision=0,
  overflow mark=arrow]{-55}{75}
\end{example}
标记检查原始值、区间端点、正负堆叠小计、对比两值和误差端点。
固定范围趋势检查所有有效观测，即使末项缺失也保留越界提示。
箭头表示数据超出量程，不表示指标好坏，也不表示超出的具体幅度。

\subsection{形状与填充}
\key{marker shape=circle|square|diamond|triangle} 对应圆、方、菱形和三角形。
\key{marker fill=solid|hollow} 控制普通点标记填充，marker size 为中心到边界的尺寸。
哑铃基准点和区间点估计固定为空心，以保留原有语义；当前点、棒棒糖和趋势点可选空心或实心。

\begin{example}[text and listing,sidebyside=false]{示例 41：四种点形状}
\gradbar[shape=lollipop,marker shape=circle]{72}

\gradbar[shape=lollipop,marker shape=square]{72}

\graddumbbell[marker shape=diamond]{48}{72}

\gradbar[shape=lollipop,marker shape=triangle,
  marker fill=hollow]{72}
\end{example}
标记形状便于黑白辨认；同一语义应保持同一形状，不应暗示不同的测量单位。

\clearpage
\section{条件配色}
\key{thresholds} 接受两个严格递增的阈值，\key{threshold colors}
接受三个颜色。颜色选择基于原始输入值，不基于四舍五入后的标签，
也不基于截断后的条长。

\begin{example}{示例 42：低、中、高三段配色}
\gradbarssetup{
  thresholds={60,80},
  threshold colors={
    lightred,gold,lightgreen
  }
}
\gradbar{59}

\gradbar{60}

\gradbar{79}

\gradbar{80}

\gradbar{92}
\end{example}

\begin{tabularx}{\linewidth}{@{}lX@{}}
\toprule
\textbf{条件} & \textbf{本例颜色}\\
\midrule
$v<60$ & lightred\\
$60\leq v<80$ & gold\\
$v\geq80$ & lightgreen\\
\bottomrule
\end{tabularx}

\begin{example}{示例 43：与目标和圆角叠加}
\gradbar[
  width=30mm,
  thresholds={60,80},
  rounded=2pt,
  target=80,
  label=end
]{72}
\end{example}

\note{启用条件配色后，所选颜色采用纯色填充，并优先于主题或
\key{negative color}。默认三色为 lightred、gold、lightgreen。
用 \key{thresholds={}} 关闭。阈值代表什么、数值越高是否越好，
应在表格标题或注释中说明。}

\clearpage
\section{指标优劣方向}
\key{better=higher} 表示越大越好，为默认值；\key{better=lower} 表示越小越好。
此参数只影响已启用的 \key{thresholds} 条件配色，不修改数值、条长、目标位置或误差范围。
没有 thresholds 时，不会自动推断评价标准，也不会改变主题配色。

\begin{example}[text and listing,sidebyside=false]{示例 44：相同数值，不同评价方向}
\gradbarssetup{max=100,width=35mm,
  thresholds={60,80},precision=0}
\begin{tabular}{@{}lcc@{}}
  \toprule
  Value & Higher is better & Lower is better\\
  \midrule
  50 & \gradbar[better=higher]{50} & \gradbar[better=lower]{50}\\
  70 & \gradbar[better=higher]{70} & \gradbar[better=lower]{70}\\
  90 & \gradbar[better=higher]{90} & \gradbar[better=lower]{90}\\
  \bottomrule
\end{tabular}
\end{example}

\key{threshold colors} 的三个颜色始终按“差、中、好”排列。
默认是浅红、金黄、浅绿。阈值 $a<b$ 始终从小到大填写：
数值区间为 $v<a$、$a\le v<b$、$v\ge b$。
higher 依次取第 1、2、3 个颜色，lower 依次取第 3、2、1 个颜色。
等于阈值时属于右侧区间；不因方向改变边界归属。

例如耗时使用 \key{thresholds={70,100},better=lower}：小于 70 为好，
70 至小于 100 为中，100 及以上为差。
堆叠组成不支持条件阈值，因此 better 不用于给各段评判优劣。
如果叠加打印预设，建议使用数值、阈值说明或纹理保留含义，避免仅靠灰度表达好坏。



\subsection{越接近目标越好、位于区间内最好}
\key{better=target} 与 \key{quality target} 配合，评价距离 $d=|v-t|$。
\key{better=interval} 与 \key{quality range={a,b}} 配合，评价到闭区间的距离
$d=\max(0,a-v,v-b)$，其中 $a\le b$。
\begin{example}[text and listing,sidebyside=false]{示例 45：目标偏差与合格范围}
\gradbar[better=target,quality target=50,
  thresholds={5,15},target=50,palette=diverging]{52}

\gradbar[better=target,quality target=50,
  thresholds={5,15},target=50,palette=diverging]{75}

\gradbar[better=interval,quality range={40,60},
  thresholds={0,10},band={40,60},palette=diverging]{55}

\gradbar[better=interval,quality range={40,60},
  thresholds={0,10},band={40,60},palette=diverging]{68}
\end{example}
这两种模式必须提供两个非负且严格递增的 thresholds，单位与原数据相同。
若 $d\le t_1$ 为好，$t_1<d\le t_2$ 为中，$d>t_2$ 为差。
\key{threshold colors} 的顺序仍为差、中、好。
例如 thresholds 为 0,10 时，区间内（含端点）才判为好，离区间不超过 10 判为中。
quality target/range 只负责评价；target/band 只负责绘制参考，需分别设置。
评价不会改变长度、端点或原始标签，也不用于堆叠、浮动区间或趋势线。

\clearpage
\section{语义配色与稳定类别纹理}
\key{palette} 管理一组颜色与纹理，与设置单条外观的 theme 和排版预设 preset 分工不同。
选项从左到右执行。配色用于表达数据角色，不能替代单位、名称与数值。

\begin{tabularx}{\linewidth}{@{}lX@{}}
\toprule
方案 & 用途\\\midrule
categorical & 六种类别色，循环使用；适合没有大小顺序的类别。\\
sequential & 五级由浅到深的蓝色，适合有序类别，不自动把数值映射到色阶。\\
diverging & 橙、浅橙、中性灰、浅青、青五色；同时设置负值、比较色及差中好颜色。\\
mono & 灰度与三种纹理组合；同时设置黑白标签和趋势线极值点色。\\
\bottomrule
\end{tabularx}
\begin{example}[text and listing,sidebyside=false]{示例 46：四套成组配色}
\gradstack[palette=categorical]{20,25,15,20,10,10}

\gradstack[palette=sequential]{20,20,20,20,20}

\gradstack[palette=diverging,min=-50]{-20,-15,10,25,30}

\gradstack[palette=mono]{35,40,25}
\end{example}
位置调色板按段序号循环。需要类别在不同顺序或子集下保持一致时，使用
\key{\gradbarscategory{类别名称}{颜色}{纹理}}，并为每条堆叠设置 stack names。
类别名称按实际文本匹配，定义遵循 TeX 分组作用域；未定义类别回退到序号配色。
\begin{example}[text and listing,sidebyside=false]{示例 47：交换类别顺序后仍保留身份}
\gradbarscategory{Compute}{gradbarsBlue}{diagonal}
\gradbarscategory{Storage}{gradbarsOrange}{dots}
\gradstack[stack names={Compute,Storage}]{60,30}

\gradstack[stack names={Storage,Compute}]{30,60}

\gradbarslegend[width=55mm]{Compute,Storage}
\end{example}
命名类别优先于序号配色，图例与数据段使用同一映射。
print 或 palette=mono 激活时，类别色被灰度方案替代，类别纹理仍保留。
纹理可选 none、diagonal、reverse、dots、crosshatch、horizontal、vertical。
类别数超过色板长度时会复用颜色；请结合名称或纹理，不能把颜色视为唯一标识。

\clearpage
\section{黑白纹理填充}
\key{print} 是一组打印预设：白色填充、浅灰轨道、黑色纹理和文字，
普通条默认使用斜线，堆叠条默认循环斜线、点阵、网格。
内部标签增加白底，以免纹理穿过文字；可通过 \key{print labels=false} 取消标签白底。

\begin{example}[text and listing,sidebyside=false]{示例 48：纹理与灰度轨道}
\gradbar[print,width=50mm,height=16pt,
  pattern=diagonal,label=auto,precision=0]{80}

\gradbar[print,width=50mm,height=16pt,
  pattern=dots,label=auto,precision=0]{60}

\gradbar[print,width=50mm,height=16pt,
  pattern=crosshatch,label=auto,precision=0]{40}
\end{example}

\key{pattern} 可选 none、diagonal、reverse、dots、crosshatch、horizontal、vertical，
分别表示无纹理、正斜线、反斜线、点阵、网格、横线和竖线。
纹理覆盖在填充色之上，保留原来的条长、圆角和截断边界。
\key{pattern color} 单独指定纹理颜色；纯色或渐变主题也可叠加纹理，不限于打印预设。

\begin{example}[text and listing,sidebyside=false]{示例 49：堆叠条与图例使用一致的纹理}
\gradstack[print,width=75mm,height=17pt,
  stack names={Train,Validation,Test},
  stack patterns={diagonal,dots,crosshatch},
  legend=true,segment labels=percent,precision=0]{60,25,15}
\end{example}

\key{stack patterns} 与段序号对应，数量不足时循环使用；零段也占用一个纹理序号。
自动图例以及 \key{\gradbarslegend} 使用相同的纹理列表。
共享图例时，让图例与条形引用同一个命名样式，以保持颜色和纹理一致。
纹理只能辅助识别，仍需保留类别名称和可读数值。打印缩放后请检查细线是否清晰。

\clearpage
\section{四套排版预设}
使用 \key{preset=paper}、\key{preset=report}、\key{preset=presentation}
或 \key{preset=outline} 一次应用一组视觉选项。
这些预设不改变量程、数据、目标值或阈值；同一选项列表仍从左到右执行。

\begin{example}[text and listing,sidebyside=false]{示例 50：四套内置排版预设}
\begin{tabular}{@{}ll@{}}
  \toprule
  Preset & Preview\\
  \midrule
  Paper & \gradbar[preset=paper,width=45mm]{82}\\
  Report & \gradbar[preset=report,width=45mm,
    target=90,thresholds={60,80}]{75}\\
  Presentation & \gradbar[preset=presentation,
    width=45mm]{92}\\
  Outline & \gradbar[preset=outline,width=45mm]{70}\\
  \bottomrule
\end{tabular}
\end{example}

\clearpage
\begin{tabularx}{\linewidth}{@{}p{30mm}X@{}}
\toprule
\textbf{预设} & \textbf{特点与设置}\\
\midrule
paper：论文简洁型 & 5pt 细条、纯色、浅灰轨道、small 字号和外部数值；标签间隔较小，上下留白各 0.5pt。\\
report：报告仪表型 & 12pt 粗条、3pt 圆角、外部数值，上下留白各 2pt；提供目标线颜色与差／中／好配色，实际目标与阈值由使用者填写。\\
\key{presentation}\newline 演示强调型 & 18pt 条形、large 字号、高对比纯色、自动黑白文字，优先使用内部标签；上下留白各 4pt。\\
outline：轮廓节墨型 & 7pt 空心轮廓、白轨道、细端点与外部数值，上下留白各 1pt；不绘制大面积填充或纹理。\\
\bottomrule
\end{tabularx}

\key{row padding} 只增加数据条自身上下的占位，不修改表格或文档的全局行距。
\key{outline=true} 单独开启空心绘制，\key{outline width} 默认 0.4pt；
outline 预设将其设为 0.35pt。堆叠条也可使用轮廓，其图例同步显示空心色块，段标签移至外部。

预设会重设主题、shape、标签位置、字体、间距与纹理等视觉设置。
例如 \key{preset=paper,shape=lollipop} 保留论文预设并改为棒棒糖；反过来写则恢复普通条。
双层条、棒棒糖、带点区间及轮廓条可能把内部标签外移，这是这些图形的避让规则。
需要打印纹理时，可在预设之后追加 print；print 会关闭轮廓绘制。

\begin{example}[text and listing,sidebyside=false]{示例 51：命名样式与新图形共用预设}
\gradbarsstyle{compact}{preset=paper,width=55mm,
  max=100,precision=0,label width=6em}
\begin{tabular}{@{}ll@{}}
  \toprule
  View & Value\\
  \midrule
  Current / reference & \gradcompare[style=compact]{65}{82}\\
  Single value & \gradbar[style=compact,shape=lollipop]{82}\\
  Lower / upper & \gradrange[style=compact]{70}{90}\\
  \bottomrule
\end{tabular}
\end{example}

\clearpage
\part{专属数据表格与数据接口}

\section{专属数据表格：gradtable}
v0.0.4 提供统一的数据表格入口。先声明列，再填写每行数据；
不需要手写对齐符号或换行命令。列类型包括 text、number、bar、bullet、interval，
以及 compare、dumbbell、stack、spark；后四种的填写方法见“更多图形列”。
\key{\gradcolumn[列选项]{标识}{标题}} 的标识必须唯一；
\key{\gradrow{单元格列表}} 按声明顺序填写逗号分隔的数据。
带逗号的文本应加花括号保护，缺失数值请显式写 NA。

\begin{example}[text and listing,sidebyside=false]{示例 52：论文型，文本、数字与普通条}
\begin{gradtable}[preset=paper,graphic width=48mm,
  text width=24mm,label width=14mm]
  \gradcolumn{method}{方法}
  \gradcolumn[type=number,options={precision=0}]{n}{样本量}
  \gradcolumn[type=bar,options={max=100,precision=1}]{score}{得分}
  \gradrow{方法 A,120,82.4}
  \gradrow{方法 B,12500,91.2}
  \gradrow{方法 C,NA,NA}
\end{gradtable}
\end{example}
列选项中的 \key{options={...}} 接受原有 gradbars 参数，如 min、max、precision、unit、
theme、target 和 references。每列共用同一组参数，单元格不单独修改尺度。
number 列只排版数值，不使用条形量程裁剪；precision、number format 和 unit 可控制显示。

\subsection{宽度和标签对齐}
graphic width 为默认图区宽度，text width 为文本列宽度，label width 为数值区的最小宽度。
同列所有外部标签预先测量，使用该列最长标签与 label width 中的较大值，统一右对齐。
标签加宽不会改变图区尺度。固定精度时小数点自然对齐，混合科学计数或单位时仅保证右对齐。
长文本会在文本列换行；默认按内容确定表宽；设置表格 width 后，可自动分配文字列和图区宽度，详见后文。
默认 graphic width=32mm、text width=22mm、label width=18mm。

\subsection{数据检查}
行数据数量必须与列数一致；列标识不能重复，所有列必须在第一个数据行或分组标题之前声明。
数字格式、区间端点和点估计的校验沿用单图接口，测量阶段的错误包含行号和列标识。
不要用省略字段表示缺失，否则列表会缺少单元格。

\clearpage
\section{更多图形列}
\subsection{双层对比条与哑铃图}
新增 \key{type=compare} 和 \key{type=dumbbell}，单元格写成“基准值/当前值”。
整格 NA 表示两者都缺失，也可使用 40/NA 或 NA/60 表示单侧缺失。
颜色、尺度、变化标注和标签选项沿用对应的独立图形命令。
输入顺序是基准/当前，默认显示标签的顺序为当前/基准；delta 显示当前减基准。

\begin{example}[text and listing,sidebyside=false]{示例 53：同一数据的两种对比表达}
\begin{gradtable}[text width=25mm,graphic width=34mm,
  label width=10mm,header align=c,column sep=7pt,
  row stretch=1.8]
  \gradcolumn{name}{方法}
  \gradcolumn[type=compare,
    options={max=100,precision=0}]{bars}{基准与当前}
  \gradcolumn[type=dumbbell,
    options={max=100,precision=0,compare label=delta,
      marker shape=diamond}]{change}{变化}
  \gradrow{方法 A,40/68,40/68}
  \gradrow{方法 B,75/60,75/60}
  \gradrow{方法 C,50/NA,50/NA}
\end{gradtable}
\end{example}
双层条适合保留两组长度，哑铃图适合强调两个端点之间的差异。
列标题应说明基准、当前或差值的含义；同一列所有行应采用同一尺度。

\clearpage
\subsection{堆叠条与迷你折线}
新增 \key{type=stack} 和 \key{type=spark}。由于行数据也使用逗号分隔，
一格内的数值列表必须用一层花括号包住，例如 \key{\gradrow{项目,{20,30,50}}}。
堆叠列显示各段及总量；迷你折线列显示序列，并在外部标注最后一个观测值。

\begin{example}[text and listing,sidebyside=false]{示例 54：构成与趋势并列展示}
\begin{gradtable}[text width=25mm,graphic width=34mm,
  label width=10mm,header align=c,column sep=7pt,
  row stretch=1.8]
  \gradcolumn{name}{项目}
  \gradcolumn[type=stack,options={max=100,precision=0,
    stack colors={gradbarsBlue,gradbarsTeal,gold}}
    ]{mix}{三类构成}
  \gradcolumn[type=spark,options={spark range=fixed,
    min=0,max=100,precision=0}]{trend}{最近五期}
  \gradrow{项目 A,{20,35,25},{35,48,NA,65,80}}
  \gradrow{项目 B,{30,20,15},{80,70,65,NA,NA}}
  \gradrow{项目 C,NA,NA}
\end{gradtable}
\end{example}
三段颜色依次代表类别一、二、三；同一类别在各行保持相同顺序和颜色。
堆叠列的整格 NA 表示构成缺失，不能把缺失的一段当作零值参与累计。
含负段时应明确设置负的 min，正段与负段分别累计。
迷你折线在缺失处断开，末次观测缺失时标签仍为缺失符号，不借用之前的值。
这里显式设置 \key{spark range=fixed,min=0,max=100}，使不同项目可按同一纵向尺度比较；
默认 auto 为各单元格独立缩放，只适合观察形态。

\clearpage
\section{表格对齐与线条样式}
\subsection{单元格对齐}
未单独设置时，表头采用对应列的对齐方式。
表格级 align=auto 时，文本和图形列默认左对齐，纯数字列默认右对齐；列选项优先于表格设置。

\begin{example}[text and listing,sidebyside=false]{示例 55：左、中、右对齐与居中表头}
\begin{gradtable}[rules=grid,header align=c,valign=c,
  text width=30mm,column sep=5pt,row padding=3pt,
  header color=gradbarsBlue!8,rule color=gradbarsBlue!60!black]
  \gradcolumn[align=l]{left}{左对齐}
  \gradcolumn[align=c]{center}{居中}
  \gradcolumn[align=r]{right}{右对齐}
  \gradcolumn[type=number,align=c,width=20mm,
    options={precision=0}]{number}{数量}
  \gradrow{方法 A,方法 A,方法 A,120}
  \gradrow{较长的名称会在文本列内自动换行,短名称,短名称,5}
\end{gradtable}
\end{example}
文本列的 \key{valign=t|c|b} 对应 array 的 p、m、b 段落列，默认 t。
“左对齐且垂直居中”可写 \key{align=l,valign=c}；
“水平和垂直均居中”可写 \key{align=c,valign=c}。
垂直设置只作用于文本列，其对齐参考遵循 p/m/b 的段落盒规则。

\clearpage
\subsection{列宽、图区宽度与数值宽度}
列选项 width 对文本列表示固定换行宽度，默认取表格 text width；
对数字和图形列表示最小单元格宽度，内容超过时自动扩展。
graphic width 只控制图区，label width 只控制数值区，三者不能相互替代。
图形列 align 对整个“图形加数值”单元格生效，图形内部数值仍保持统一右对齐。

\begin{example}[text and listing,sidebyside=false]{示例 56：得分列与表头居中}
\begin{gradtable}[rules=rows,header align=c,
  graphic width=28mm,label width=12mm,
  column sep=8pt,row stretch=1.4]
  \gradcolumn[width=24mm,header align=l]{method}{方法}
  \gradcolumn[type=bar,align=c,width=75mm,
    options={max=100,precision=0}]{score}{得分}
  \gradrow{方法 A,82}
  \gradrow{方法 B,NA}
\end{gradtable}
\end{example}

\subsection{四种线条样式和间距参数}
线条与图形主题独立，可与 paper、report、mono 任意组合。

\begin{tabularx}{\linewidth}{@{}lX@{}}
\toprule
选项 & 含义与默认值\\\midrule
\key{rules=booktabs} & 默认三线表：顶线、表头分隔线、底线。\\
\key{rules=rows} & 三线表基础上增加数据行之间的横线。\\
\key{rules=grid} & 完整横纵网格与外边框。\\
\key{rules=none} & 不绘制表格线条。\\
\key{column sep} & 每列左右内边距，默认 6pt；相邻列通常合计为其两倍。\\
\key{row stretch} & 基础行高系数；默认 paper/mono 为 1.2，report 为 1.5。\\
\key{row padding} & 附加行高，默认 0pt；不缩放单元格图形。\\
\key{header color} & 表头底色；默认 report 为 gradbarsTrack，其余无底色，空值清除。\\
\key{rule color} & 表格线颜色，默认 black，设置只影响当前 gradtable。\\
\bottomrule
\end{tabularx}

行高系数必须为正，列间距、附加行高及列宽不能为负。
width=0pt 表示使用默认：文本列采用 text width，其余列按内容排版。

\begin{example}[text and listing,sidebyside=false]{示例 57：无线框、统一居中与单列覆盖}
\begin{gradtable}[rules=none,align=c,header align=c,
  text width=30mm,column sep=10pt,row stretch=1.5,
  header color=gradbarsTeal!8]
  \gradcolumn{name}{项目}
  \gradcolumn[type=number,width=24mm,
    options={precision=0}]{count}{数量}
  \gradcolumn[type=number,align=r,width=24mm,
    options={precision=1,unit={ 分}}]{score}{评分}
  \gradrow{项目 A,12,82.4}
  \gradrow{项目 B,120,91.2}
\end{gradtable}
\end{example}
这些 align 参数控制单元格，不负责整张表在页面上的位置。
整张表居中时，可在外层 table 环境中使用标准的 \key{\centering}；
左、右放置可使用外层分组中的 \key{\raggedright} 或 \key{\raggedleft}。

\clearpage
\section{按总宽度分配列宽}
表格选项 \key{width=\linewidth} 将整张表限定为当前可用宽度。
宏包先测量该列全部数值标签，保留标签、间距与线条的空间，
再为文字列和图形区域分配剩余宽度。图形的数值尺度保持不变，字体不会被整体缩小。

\begin{example}[text and listing,sidebyside=false]{示例 58：固定数值列与可伸缩图形列}
\begin{gradtable}[width=\linewidth,header align=c,
  label width=12mm,column sep=6pt,row padding=2pt,
  min text width=18mm,min graphic width=15mm]
  \gradcolumn[weight=1]{method}{方法名称}
  \gradcolumn[type=bar,weight=2,mark=both,
    options={max=100,precision=1}]{score}{得分}
  \gradcolumn[type=number,width=18mm,
    options={precision=0}]{count}{样本量}
  \gradheader{\gradspan{1}{实验设置}
    \gradspan{2}{评估结果}}
  \gradrow{带有较长名称的改进方法,92.8,12500}
  \gradrow{基础方法,84.6,320}
  \gradrow{待测方法,NA,80}
\end{gradtable}
\end{example}
\begin{tabularx}{\linewidth}{@{}lX@{}}
\toprule
参数 & 分配规则\\\midrule
\key{width=0pt} & 默认值，保持按内容排版；正值启用总宽度分配。\\
\key{weight} & 列选项，默认 1。无固定宽度的文字、图形列按权重分配剩余空间。\\
\key{min text width} & 可伸缩文字列的最低宽度，默认 12mm。\\
\key{min graphic width} & 可伸缩图区的最低宽度，默认 12mm，不含外部数值。\\
列选项 \key{width} & 总宽度模式下固定文字或图形列的单元格宽度；数字列仍至少容纳全部标签。\\
\bottomrule
\end{tabularx}
权重作用于满足最低宽度之后的剩余空间，因此 weight=2 不代表最终列宽恰为另一列的两倍。
数字列按标签所需宽度排版，不参与伸缩。图形列的最终图区宽度由分配结果决定，
优先于 graphic width 或图形 options 内的 width；同一列的所有行始终共享同一图区宽度。

\note{总宽度不足以容纳固定列、标签、最小图区及间距时，宏包会报错。
可以增加表宽、减小 column sep 或最低宽度、缩短单位，或减少列数。
所有列都固定后若仍有剩余空间，应至少保留一个可伸缩文字或图形列。
总宽度模式允许长表头和分组标题换行；文字中的不可断长串仍需自行提供断行位置。}

\clearpage
\section{多级表头与跨列标题}
\key{\gradheader{...}} 添加一层表头，内部用 \key{\gradspan{列数}{标题}} 依次声明合并区域。
多次调用 gradheader 时按声明顺序从上到下排列，最下层仍显示各 gradcolumn 的标题。
先声明全部列，再声明表头，最后填写分组和数据。每层合并列数之和必须恰好等于总列数。

\begin{example}[text and listing,sidebyside=false]{示例 59：三级表头组织不同指标}
\begin{gradtable}[rules=grid,header align=c,
  header color=gradbarsBlue!8,text width=28mm,
  column sep=8pt,row padding=3pt]
  \gradcolumn{method}{方法}
  \gradcolumn[type=number,width=24mm,
    options={precision=1}]{accuracy}{准确率}
  \gradcolumn[type=number,width=24mm,
    options={precision=1}]{fscore}{F1}
  \gradcolumn[type=number,width=24mm,
    options={precision=0}]{latency}{毫秒}
  \gradheader{\gradspan{4}{模型评估}}
  \gradheader{\gradspan{1}{方法信息}
    \gradspan{2}{预测质量}\gradspan{1}{推理耗时}}
  \gradrow{方法 A,91.4,89.7,24}
  \gradrow{方法 B,93.2,91.5,31}
\end{gradtable}
\end{example}
跨列标题默认居中，底色跟随 header color，分隔线跟随 rules。
空标题可写成 \key{\gradspan{1}{}}，该单元格仍占指定列数。
本接口合并同一行中的连续列；各层标题独立排布，不执行跨行合并。
命名列定义中也可保存 gradheader，但在表头之后不能再追加列声明。

\clearpage
\section{分组行与最优值标记}
\subsection{用分组标题组织数据}
\key{\gradgroup{标题}} 插入跨越所有列的分组行。先声明全部列，再填写分组和数据；
每个分组必须至少有一条数据行。分组不占数据列，不参与数值比较，也不会改变列尺度。
\key{group color} 设置标题底色，默认 gradbarsTrack，空值表示无底色；
\key{group sep} 默认 4pt，为分组标题增加上方留白，非网格表还增加组前间距。
分组行兼容四种 rules 样式。标题过长时应缩短文字，避免把整张表撑宽。

\begin{example}[text and listing,sidebyside=false]{示例 60：分组展示与全表最优值}
\begin{gradtable}[header align=c,text width=30mm,
  graphic width=38mm,label width=12mm,group sep=4pt]
  \gradcolumn{method}{方法}
  \gradcolumn[type=bar,mark=both,
    options={max=100,precision=1}]{score}{得分}
  \gradcolumn[type=number,mark=both,width=24mm,
    options={better=lower,precision=1}]{time}{耗时 / 秒}
  \gradgroup{基础方法}
  \gradrow{方法 A,82.4,12.8}
  \gradrow{方法 B,90.2,10.5}
  \gradgroup{改进方法}
  \gradrow{方法 C,90.2,9.6}
  \gradrow{方法 D,86.7,9.6}
  \gradrow{方法 E,NA,NA}
\end{gradtable}
\end{example}
\textbf{最佳值加粗，次佳值加下划线。} 得分的两个 90.2 均加粗，86.7 为次佳值；
耗时的两个 9.6 均加粗，10.5 为次佳值。比较覆盖当前整张表的全部数据行，跨越分组。

\clearpage
\subsection{排名规则与数据含义}
列选项 \key{mark=none|best|both} 默认为 none；best 仅标记最佳值，both 同时标记次佳值。
该选项适用于 number、bar 和 bullet 列，只改变数字标签，不改变条长或原始数据。
在列的 options 内使用 \key{better=higher}（默认）或 \key{better=lower} 指定方向；
目标接近度和区间达标不能直接用于本接口的最优值排序。

\begin{example}[text and listing,sidebyside=false]{示例 61：负数、并列值与零值}
\begin{gradtable}[rules=grid,header align=c,
  text width=34mm,column sep=8pt,row padding=2pt]
  \gradcolumn{method}{项目}
  \gradcolumn[type=number,mark=both,width=28mm,
    options={better=lower,precision=1}]{loss}{偏差}
  \gradrow{方案 A,-2}
  \gradrow{方案 B,-2.0}
  \gradrow{方案 C,0}
  \gradrow{方案 D,NA}
\end{gradtable}
\end{example}
\begin{itemize}
  \item 排名采用格式化之前的原始数值；不会按字符串或四舍五入后的标签排序。
  \item 数值相等即为并列，包括 2、2.0 等写法；次佳值取下一个不同的数值。
  \item NA 不参与排名，零值正常参与。整列缺失时没有标记，全部相等时没有次佳值。
  \item 若精度过低，不同原始数值可能显示成相同标签；应提高 precision，避免标记难以解释。
  \item 使用加粗或下划线前，应在表注中解释约定。堆叠、趋势和对比列没有唯一的默认排名指标，
        请另设 number 列明确指标，再标记最优值。
\end{itemize}

\clearpage
\section{报告与黑白表格主题}
表格 preset 可选 \key{paper}（三线表、纯色细条、紧凑行距）、
\key{report}（浅色表头、较高条形、宽行距）或 \key{mono}（灰度和纹理）。
主题先应用，列 options 随后覆盖。表格主题与单图预设分别设置，单图仍可用于普通 tabular。

\begin{example}[text and listing,sidebyside=false]{示例 62：报告型，子弹图与区间列}
\begin{gradtable}[preset=report,graphic width=32mm,
  text width=19mm,label width=19mm]
  \gradcolumn{name}{项目}
  \gradcolumn[type=bullet,options={max=100,precision=0,
    bullet bands={{60/black!8},{85/black!16},{100/black!25}},
    target=85}]{actual}{完成度}
  \gradcolumn[type=interval,options={max=100,precision=0,
    color=gradbarsTeal}]{estimate}{估计区间}
  \gradrow{项目 A,78,65/78/90}
  \gradrow{项目 B,92,80/92/98}
  \gradrow{项目 C,NA,NA}
\end{gradtable}
\end{example}
interval 列使用 \key{下界/点估计/上界}，点估计为 NA 时仅保留区间。
整个字段为 NA 表示区间缺失；点估计必须位于原始端点之间。
子弹图必须显式提供背景等级，本例 target 是完成目标，不表示区间点估计。

\begin{example}[text and listing,sidebyside=false]{示例 63：黑白型，正负变化与区间}
\begin{gradtable}[preset=mono,graphic width=32mm,
  text width=19mm,label width=19mm]
  \gradcolumn{name}{组别}
  \gradcolumn[type=bar,options={min=-40,max=60,precision=0}]{change}{变化量}
  \gradcolumn[type=interval,options={min=-40,max=60,
    precision=0,marker shape=diamond}]{range}{区间}
  \gradrow{组 A,-12,-25/-12/5}
  \gradrow{组 B,28,15/28/45}
\end{gradtable}
\end{example}
默认 gradtable 不分页，可放入 table、table* 或 minipage；long=true 模式直接用于单栏正文。
表格不自动合计、推断量程或读取数据文件。
大量行可使用 gradtable 的 long=true 模式，或继续使用 longtable 与单图命令。

\clearpage
\section{复用表格样式与列定义}
\key{\gradtablestyle{名称}{表格选项}} 保存整张表的外观，
\key{\gradtableschema{名称}{列声明}} 保存列顺序、标题、类型和列参数。
在环境选项中使用 \key{style=名称}，在环境内部使用 \key{\graduseschema{名称}}。
定义一般放在导言区；重复定义同名项目会替换之前的定义。

\begin{example}[text and listing,sidebyside=false]{示例 64：复用同一套样式与列定义}
\gradtablestyle{benchmark}{rules=booktabs,header align=c,
  text width=30mm,graphic width=40mm,label width=12mm,
  row stretch=1.3,group color=gradbarsBlue!8}
\gradtableschema{benchmark}{
  \gradcolumn{method}{方法}
  \gradcolumn[type=bar,mark=both,
    options={max=100,precision=1}]{score}{得分}
}
\begin{gradtable}[style=benchmark]
  \graduseschema{benchmark}
  \gradgroup{验证集}
  \gradrow{方法 A,81.5}
  \gradrow{方法 B,88.2}
\end{gradtable}

\begin{gradtable}[style=benchmark,rules=grid]
  \graduseschema{benchmark}
  \gradgroup{测试集}
  \gradrow{方法 A,79.6}
  \gradrow{方法 B,85.3}
\end{gradtable}
\end{example}
选项按从左到右的顺序执行；放在 style 后面的参数覆盖命名样式。
样式可以引用其他已定义样式，列定义也可以组合其他列定义；循环引用会报告错误。
列标识符在一张表内必须唯一，列定义必须在第一个数据行或分组标题之前调用。
每张表分别计算列宽和最优值，不会把前一张表的数据带入下一张表。
默认 gradtable 不分页；设置 long=true 可启用跨页模式，详见后文。当前不自动计算汇总行。

\clearpage
\section{跨页表格与重复表头}
\key{long=true} 使用 longtable 排版，允许在数据行之间分页。
完整的多级表头会在每一页重复；列宽与最优值在整张表中统一计算，不会逐页改变。
分组标题与首条数据保持在同一页，后续数据可以跨页，续页不自动重复分组名称。

\begin{tabularx}{\linewidth}{@{}lX@{}}
\toprule
选项 & 用法\\\midrule
\key{long=true} & 开启跨页；默认 false。直接放在单栏正文中。\\
\key{caption={标题}} & 长表标题；自动编号，续页沿用同一编号。\\
\key{label={tab:results}} & 用于标准 ref 引用；须同时提供 caption。\\
\key{continued text={（续）}} & 续页标题及非末页页脚提示，默认 (continued)，可设为空。\\
\bottomrule
\end{tabularx}
长表不能放在 table 浮动体、minipage、双栏模式或不可分页的示例框中。
单条数据行本身不能拆分；一个超高单元格应先拆成多条记录。
跨页表默认靠左放置，width 可设为当前 linewidth。页数和引用通常需要编译两次才稳定。

\subsection{示例 65：多级表头、自动列宽与跨页的组合}
下表直接在手册正文中分页，随后给出产生它的完整源码。
图形和标签尺度跨页保持一致；加粗及下划线比较的是全部记录。
\begin{filecontents*}[overwrite]{\jobname-gradtable-long.tex}
\begin{gradtable}[long=true,width=\linewidth,
  caption={分组模型评估},label={tab:gradtable-long},
  continued text={（续）},header align=c,
  label width=12mm,column sep=6pt,
  row stretch=1.4,row padding=3pt]
  \gradcolumn[weight=1]{method}{方法}
  \gradcolumn[type=bar,weight=2,mark=both,
    options={max=100,precision=1}]{score}{得分}
  \gradcolumn[type=number,width=18mm,mark=both,
    options={better=lower,precision=1}]{time}{耗时 / 秒}
  \gradcolumn[type=interval,weight=1.5,
    options={max=100,precision=0}]{range}{区间}
  \gradheader{\gradspan{1}{实验设置}
    \gradspan{3}{同一测试集的评估结果}}
  \gradgroup{实验组 A}
  \gradrow{方法 01,70.0,8.0,66/70/74}
  \gradrow{方法 02,77.2,11.1,73/77/81}
  \gradrow{方法 03,84.4,14.2,80/84/88}
  \gradrow{方法 04,91.0,17.0,87/91/95}
  \gradrow{方法 05,70.2,20.1,66/70/74}
  \gradrow{方法 06,77.4,23.2,73/77/81}
  \gradrow{方法 07,84.0,10.0,80/84/88}
  \gradrow{方法 08,91.2,13.1,87/91/95}
  \gradgroup{实验组 B}
  \gradrow{方法 09,70.0,16.0,66/70/74}
  \gradrow{方法 10,NA,19.1,NA}
  \gradrow{方法 11,84.4,22.2,80/84/88}
  \gradrow{方法 12,91.0,9.0,87/91/95}
  \gradrow{方法 13,70.2,12.1,66/70/74}
  \gradrow{方法 14,77.4,15.2,73/77/81}
  \gradrow{方法 15,84.0,18.0,80/84/88}
  \gradrow{方法 16,91.2,21.1,87/91/95}
  \gradgroup{实验组 C}
  \gradrow{方法 17,70.0,8.0,66/70/74}
  \gradrow{方法 18,77.2,11.1,73/77/81}
  \gradrow{方法 19,84.4,14.2,80/84/88}
  \gradrow{方法 20,91.0,17.0,87/91/95}
  \gradrow{方法 21,70.2,20.1,66/70/74}
  \gradrow{方法 22,77.4,23.2,73/77/81}
  \gradrow{方法 23,84.0,10.0,80/84/88}
  \gradrow{方法 24,91.2,13.1,87/91/95}
\end{gradtable}
\end{filecontents*}
\input{\jobname-gradtable-long.tex}
表~\ref{tab:gradtable-long} 的数据列在续页沿用相同的位置与宽度。

\clearpage
\subsection{示例 65 的完整源码}
以下源码与上面的实际长表来自同一个内容块，修改数据后再次编译即可更新。
\lstinputlisting[style=manual,backgroundcolor=\color{manualCode},
  basicstyle=\ttfamily\fontsize{8}{10.2}\selectfont]{\jobname-gradtable-long.tex}

\clearpage
\section{CSV 数据导入}
\key{\gradbarsloadcsv[选项]{数据集名}{文件名}} 将 CSV 读入当前 TeX 作用域，
随后可生成数据表、逐行取值，或从指定数值列计算共享量程。
文件以 UTF-8 保存，首行是列名；列名区分大小写。
以下所有命令从仓库根目录编译，使用配套文件 \key{docs/data/results.csv}。

\textbf{示例数据文件}\par
\lstinputlisting[style=manual,language={}]{docs/data/results.csv}

\begin{example}[text and listing,sidebyside=false]{示例 66：从 CSV 直接生成数据条表格}
\gradbarsloadcsv{results}{docs/data/results.csv}
\gradbarscsvtable[width=40mm,precision=1]{results}{Model}{Score}
\end{example}

表格按文件顺序显示 Model 与 Score 两列。数值列自动采用同一范围：
下限为零和有效值最小值的较小者，上限为零和有效值最大值的较大者。
若上限为零，则使用 1，保证全零、全缺失或全负数列仍有有效量程。
空值与缺失标记不参与范围计算。

默认缺失标记为 \key{NA}、\key{N/A}、\key{null}，区分大小写；空字段始终表示缺失。
加载时可用 \key{missing={NA,N/A,null,--}} 替换缺失标记列表。
映射仅作用于数值读取，文本列中的相同字符串保留原样。
零值仍是有效观测，不能用缺失标记替代。

\note{CSV 支持逗号分隔、双引号包裹字段、字段内逗号、用两个双引号表示一个引号，
以及文件开头的 UTF-8 BOM。字段首尾空白会被去除，空白行会跳过。
不支持字段内换行、分号或制表符分隔；未闭合引号、重复或空列名、行列数不匹配会报错。
文件内容按普通字符读取，不执行其中的 TeX 命令。}

\clearpage
\subsection{按列建立共享样式，保留自己的表格布局}
\key{\gradbarscsvstyle{样式名}{数据集名}{列名}} 扫描指定列，
生成包含 min 和 max 的命名样式。生成样式会覆盖同名旧样式，
使用时可继续追加宽度、单位等显示选项。
\key{\gradbarcsv[选项]{数据集名}{行号}{列名}} 绘制单个数据条，行号从 1 开始，不包含表头。
\key{\gradbarscsvcell{数据集名}{行号}{列名}} 显示原始文本字段。

\begin{example}[text and listing,sidebyside=false]{示例 67：共享耗时范围与自定义表格}
\gradbarsloadcsv{results}{docs/data/results.csv}
\gradbarscsvstyle{csvLatency}{results}{Latency}
\gradbarssetup{style=csvLatency,width=40mm,
  unit={\,ms},precision=0}
\begin{tabular}{@{}ll@{}}
  \toprule
  Model & Latency\\
  \midrule
  \gradbarscsvcell{results}{1}{Model} &
    \gradbarcsv{results}{1}{Latency}\\
  \gradbarscsvcell{results}{3}{Model} &
    \gradbarcsv{results}{3}{Latency}\\
  \gradbarscsvcell{results}{4}{Model} &
    \gradbarcsv{results}{4}{Latency}\\
  \bottomrule
\end{tabular}
\end{example}

本例共享范围为 0--125。自动范围包含整个数值列，而不是仅包含表中选取的几行。
修改 CSV 后重新编译，数据与范围都会更新。若要跨批次比较，建议明确覆盖固定的
\key{min} 和 \key{max}，避免每批数据的最大值变化导致视觉尺度漂移。

\key{\gradbarscsvtable} 生成普通双列表格，适合简短展示；不会自动分页。
需要长表、多列或复杂表头时，保留自己的 longtable/tabular 布局，通过取值命令逐行填入。
数值必须是十进制字面量，不支持千位分隔、科学计数输入、公式或单位后缀。
其他非数字内容会报错，不会静默丢弃。数据集名使用英文字母开头，可包含字母、数字、下划线和连字符。

\clearpage
\part{完整排版案例}

\section{完整案例：模型对比}
这张表使用两个独立的列范围：准确率最大值为 100，延迟最大值为 50。
条越长只表示数值越大；准确率通常越高越好，延迟通常越低越好。
在下面的宽幅案例中，效果放在上方、源码放在下方，避免缩小表格字体。

\begin{example}[text and listing,sidebyside=false]{示例 68：准确率与延迟}
\begingroup
\gradbarssetup{
  width=27mm, height=1.5ex,
  label width=3.6em, precision=1
}
\renewcommand{\arraystretch}{1.65}
\begin{tabular}{@{}lcc@{}}
  \toprule
  Model & Accuracy (\%) & Latency (ms) \\
  \midrule
  Baseline
    & \gradbar{72.4}
    & \gradbar[theme=teal,max=50]{18.2} \\
  Compact
    & \gradbar{81.6}
    & \gradbar[theme=teal,max=50]{12.5} \\
  Transformer
    & \gradbar{89.3}
    & \gradbar[theme=teal,max=50]{42.8} \\
  \textbf{Proposed}
    & \gradbar{92.7}
    & \gradbar[theme=teal,max=50]{24.1} \\
  \bottomrule
\end{tabular}
\endgroup
\end{example}

\subsection{嵌入自己的论文}
在主文档导言区加载 \key{booktabs} 和 \key{gradbars}，然后复制上述代码即可。
需要标题和交叉引用时，再在外层使用普通 \key{table} 环境、
\key{\caption} 和 \key{\label}。
同一列应保持相同的最大值、轨道宽度和标签槽宽。

\clearpage
\section{完整案例：数据集概览}
下面将同一组数据分别显示为样本数和总量占比。两列分母都为 125000；
标签格式不同，但相同数据的彩色条长度保持一致。

\begin{example}[text and listing,sidebyside=false]{示例 69：样本量与占比}
\begingroup
\gradbarssetup{
  max=125000, width=29mm,
  label width=4.8em, precision=0
}
\renewcommand{\arraystretch}{1.65}
\begin{tabular}{@{}lcc@{}}
  \toprule
  Split & Samples & Share of total \\
  \midrule
  Training
    & \gradbar[theme=solid]{87500}
    & \gradbar[theme=teal,
        value format=percent,precision=1]{87500} \\
  Validation
    & \gradbar[theme=solid]{18750}
    & \gradbar[theme=teal,
        value format=percent,precision=1]{18750} \\
  Test
    & \gradbar[theme=solid]{18750}
    & \gradbar[theme=teal,
        value format=percent,precision=1]{18750} \\
  Unassigned
    & \gradbar[theme=solid]{0}
    & \gradbar[theme=teal,
        value format=percent,precision=1]{0} \\
  \bottomrule
\end{tabular}
\endgroup
\end{example}

\note{本例演示较大数值与真实零值。样本数中的分组空白仅用于可读性，
输入仍必须写成 \key{87500}，不能在数值参数中加入逗号或单位。}

\clearpage
\section{完整案例：主题与标签总览}
这个案例把四种主题、三种标签位置、自定义标签和正文用法放在一起，
可直接复制到文档中比较效果。

\begin{example}[text and listing,sidebyside=false]{示例 70：主题与布局速查}
\begingroup
\gradbarssetup{
  width=30mm, label width=3.8em,
  unit={\%}, precision=0
}
\renewcommand{\arraystretch}{1.6}
\begin{tabular}{@{}ll@{}}
  \toprule
  Style & Preview \\
  \midrule
  Blue    & \gradbar[theme=blue]{72} \\
  Teal    & \gradbar[theme=teal]{72} \\
  Solid   & \gradbar[theme=solid]{72} \\
  Gray    & \gradbar[theme=gray]{72} \\
  Inside  & \gradbar[label=inside,
              height=2.5ex,right color=white]{64} \\
  Custom  & \gradbar[theme=teal,max=10,
              text={8 / 10}]{8} \\
  No label & \gradbar[theme=gray,label=none]{64} \\
  \bottomrule
\end{tabular}

Completion:
\gradbar[width=16mm,label width=3em]{64}.
\endgroup
\end{example}

\subsection{如何选用}
为一组指标选择稳定的样式即可，不必为每个数值选择不同的颜色。
需要黑白输出时选择灰度；需要减少视觉装饰时选择纯色。
精确比较依靠数值标签，彩色条提供快速浏览的线索。

\clearpage
\section{完整案例：前后对比与目标评价}
下面用同一尺度比较方法改进前后的得分，并将目标偏差单独表达。
数据为演示值；哑铃图例写在表注，避免读者混淆两个端点。
\begin{example}[text and listing,sidebyside=false]{示例 71：实验前后对比表}
\gradbarsstyle{study}{preset=paper,width=48mm,
  min=0,max=100,precision=0,palette=categorical}
\begin{tabular}{@{}lll@{}}
\toprule
Method & After / Before & Target 80\\\midrule
A & \graddumbbell[style=study]{58}{76}
  & \gradbar[style=study,width=23mm,better=target,
    quality target=80,thresholds={5,15}]{76}\\
B & \graddumbbell[style=study]{63}{84}
  & \gradbar[style=study,width=23mm,better=target,
    quality target=80,thresholds={5,15}]{84}\\
C & \graddumbbell[style=study]{65}{65}
  & \gradbar[style=study,width=23mm,better=target,
    quality target=80,thresholds={5,15}]{65}\\
\bottomrule
\end{tabular}
\end{example}
空心圆为 Before，实心圆为 After；目标评价条越长仍表示数值越大。
颜色表示到 80 的偏差，而非“越长越好”。

\clearpage
\section{完整案例：收支贡献与黑白打印}
正负分开累计，同一类别跨行保持纹理一致；表头明确显示的是正负小计。
这不是瀑布图：每一段不是接着上一段净值往上或往下移动。
\begin{example}[text and listing,sidebyside=false]{示例 72：共享类别图例的贡献表}
\gradbarscategory{Revenue}{gradbarsBlue}{diagonal}
\gradbarscategory{Materials}{gradbarsOrange}{dots}
\gradbarscategory{Service}{gradbarsGreen}{crosshatch}
\gradbarsstyle{ledger}{width=65mm,height=10pt,
  min=-60,max=100,precision=0,palette=mono,
  stack names={Revenue,Materials,Service},
  stack totals=separate}
\begin{tabular}{@{}ll@{}}
\toprule
Period & Positive / Negative (units)\\\midrule
Q1 & \gradstack[style=ledger]{80,-35,-10}\\
Q2 & \gradstack[style=ledger]{95,-30,-15}\\
Q3 & \gradstack[style=ledger]{70,-40,-20}\\
\bottomrule
\end{tabular}

\gradbarslegend[palette=mono,width=65mm]{Revenue,Materials,Service}
\end{example}
所有行都使用 $[-60,100]$ 范围，零点位置固定。图例在表外只出现一次。

\clearpage
\section{完整案例：等间隔趋势汇总}
同表采用固定纵向范围，不逐行放大微小波动；缺失月份保留位置。
\begin{example}[text and listing,sidebyside=false]{示例 73：半年观测与末期数值}
\gradbarsstyle{monthly}{width=65mm,height=7mm,
  min=0,max=100,spark range=fixed,precision=0,
  row padding=3pt,spark points=extrema}
\begin{tabular}{@{}ll@{}}
\toprule
Series & Jan--Jun / June value\\\midrule
North & \gradspark[style=monthly]{35,48,42,60,72,80}\\
South & \gradspark[style=monthly]{55,53,NA,58,62,65}\\
West & \gradspark[style=monthly]{45,45,45,45,45,45}\\
East & \gradspark[style=monthly]{30,40,55,60,70,NA}\\
\bottomrule
\end{tabular}
\end{example}
North 与 South 的极值点只表示各自的最高与最低观测，不表示好坏。
East 的末期标签为缺失，不能用前一期的 70 代替。
不等间隔日期应先整理到公共观测网格；不要直接当作等间距序列。

\clearpage
\section{真实排版案例}
下面四个案例使用各自的文档类实际编译，效果与完整源码均收录于本手册。
配套源码位于 \key{docs/layouts}；从仓库根目录编译可找到 \key{gradbars.sty}。
这些是实际排版环境，所有数值仍为演示数据。

\subsection{示例 74：双栏论文}
使用 article 的 twocolumn 选项；数据条宽度按单栏分配。
下图是完整双栏页面。窄列的总宽度还需加上外部标签槽与间隔。
\begin{center}
\fbox{\includegraphics[width=.84\linewidth,height=.72\textheight,keepaspectratio]{docs/layouts/twocolumn.pdf}}
\end{center}
\clearpage
\textbf{示例 74 完整源码}\par
\lstinputlisting[style=manual]{docs/layouts/twocolumn.tex}

\clearpage
\subsection{示例 75：跨页长表}
longtable 在单栏文档中自动分页，并在后续页重复表头。
该案例包含 60 行，实际输出两页。所有页面使用相同的列定义和 0--100 量程。
\begin{center}
\fbox{\includegraphics[page=1,width=.46\linewidth]{docs/layouts/longtable.pdf}}
\hfill
\fbox{\includegraphics[page=2,width=.46\linewidth]{docs/layouts/longtable.pdf}}
\end{center}
长表不能放在 table 浮动体内，也不能直接用在 article 的双栏模式中。
需要跨页时，将这部分切换为单栏或单独编排。
\clearpage
\textbf{示例 75 完整源码}\par
\lstinputlisting[style=manual,basicstyle=\ttfamily\fontsize{7.5}{9}\selectfont]{docs/layouts/longtable.tex}

\clearpage
\subsection{示例 76：黑白打印}
使用灰度主题与不同灰度的分段，同时保留数值、图例名称和目标线。
下图截取实际页面的内容区域，堆叠条与图例同时使用黑白纹理。
\begin{center}
\fbox{\includegraphics[trim=35bp 360bp 35bp 35bp,clip,width=.72\linewidth]{docs/layouts/grayscale.pdf}}
\end{center}
\textbf{示例 76 完整源码}\par
\lstinputlisting[style=manual]{docs/layouts/grayscale.tex}

\clearpage
\subsection{示例 77：幻灯片}
使用 Beamer 的 16:9 页面。增大轨道高度，并通过 arraystretch 增加表格行距，
使标签和相邻条形之间保留空隙。若幻灯片内还要显示 verbatim 源码，需给 frame 添加 fragile。
\begin{center}
\fbox{\includegraphics[width=.92\linewidth]{docs/layouts/slides.pdf}}
\end{center}
\textbf{示例 77 完整源码}\par
\lstinputlisting[style=manual,basicstyle=\ttfamily\fontsize{7.5}{8.5}\selectfont]{docs/layouts/slides.tex}

\clearpage
\part{常见问题与接口参考}

\section{常见问题}
\subsection{长标签仍然超出页面怎么办？}
\key{label overflow=auto} 防止标签挤在轨道中，但不能突破页面宽度限制。
外部槽宽会按文字扩展，因此应缩短文字、减小轨道宽度，或把长说明移到相邻文本列。
需要多行说明时，可在 \key{text} 中提供指定宽度的 \key{\parbox}，并预留行高。
若整列标签必须右对齐，请将 \key{label width} 设为能容纳该列最长标签的统一宽度。

\subsection{条形相贴，或行高太大怎么办？}
轨道高度由 \key{height} 设置；表格行距可通过
\key{\renewcommand{\arraystretch}{1.5}} 或 array 的
\key{\setlength{\extrarowheight}{3pt}} 调整。
自动外移标签和图例会增加图形高度，这是为文字保留实际空间。
密集表格可把自动图例关闭，表外只放一个 \key{\gradbarslegend}。

\subsection{同列的条形为什么不可直接比较？}
请检查各行的 \key{min}、\key{max}、\key{width} 是否一致。
正负条尤其需要同时统一上下限；仅统一最大值仍可能让零点位置不同。
建议用一个命名样式或 \key{\gradbarscolumn} 定义比较列。
手写数值列不会根据当前行的值推断范围；CSV 接口可根据整列计算共享范围。

\subsection{缺失值和零值有何区别？}
空参数、空数值单元格以及大写 \key{NA} 表示缺失；默认显示 \key{--}，
可用 \key{missing text} 修改。\key{0} 是有效观测，显示零标签与空轨道。
缺失值不应替换为零。堆叠组成不接受缺失段，因为未知段会使总量和占比不明确。

\subsection{数值列为什么报 Invalid value？}
G 列只接收十进制数、能展开为十进制数的宏、空值或 \key{NA}。
不要在单元格里写百分号、单位或普通表头；单位放到列选项，表头用
\key{\multicolumn{1}{c}{Score}}。需要手写特殊内容时同样用 multicolumn 绕过收集。

\subsection{标签占比为什么与条长不相同？}
普通条总标签的 \key{value format=percent} 使用量程上限作为分母
（仅支持 \key{min=0}）。\par
堆叠段的 \key{segment labels=percent} 使用分段总和作为分母。
二者回答的问题不同；示例 17 将总量占比和组成占比并列展示。

\subsection{支持哪些表格和编译方式？}
本手册验证 XeLaTeX 下的 tabular、booktabs、longtable 和 Beamer 场景。
tabularray 的专有列处理、siunitx 的 S 列及其他引擎不在当前验证范围。
正文中的普通数据条与表格中的数据条使用相同命令。



\subsection{正负堆叠的净值为零，为什么仍然有条形？}
正负贡献分别累计，50 与 $-50$ 的净值为零但活动量为 100。
需要看到两侧小计时使用 stack totals=separate；分段百分比的分母为绝对值之和。
\subsection{两条趋势线看起来一样高，数值却不同？}
默认 auto 范围只比较形状。要比较绝对变化，请共用 spark range=fixed、min、max 和 height。
\subsection{重排类别后颜色改变怎么办？}
位置色板按序号取色。定义 gradbarscategory 并提供 stack names，可保持类别身份；图例使用相同名称。

\clearpage
\section{参数速查与使用限制}
\begingroup\small
\begin{tabularx}{\linewidth}{@{}p{31mm}p{25mm}X@{}}
\toprule
\textbf{参数} & \textbf{默认值} & \textbf{含义与可选值}\\
\midrule
\key{min} & \key{0} & 范围下限，必须为零或负值。\\
\key{max} & \key{100} & 范围上限，必须为正值。\\
\key{width} & \key{24mm} & 正的轨道宽度，不含外部标签。\\
\key{height} & \key{1.5ex} & 正的轨道高度。\\
\key{theme} & \key{blue} & blue、teal、solid、gray、lbyellow、viblue、cyblu。\\
\key{label} & \key{outside} & outside、inside、none、end、auto。\\
\key{label width} & \key{4.5em} & 外部标签最小槽宽；自动溢出模式下可扩展。\\
\key{label gap} & \key{.6em} & 轨道与外部标签之间的间隔。\\
\key{precision} & \key{1} & 0--6 的整数；保留末尾零。\\
\key{value format} & \key{value} & 原始数值 value 或计算占比 percent。\\
\key{unit} & 空 & 原始数值的后缀；占比模式忽略此项。\\
\key{text} & 空 & 覆盖整个标签；空值恢复自动格式。\\
\key{overflow} & \key{clip} & clip 警告并截断，或 error 报错。\\
\key{color} & 未指定 & 指定单色，并切换为纯色填充。\\
\key{fill} & 由主题指定 & gradient 或 solid。\\
\key{left color} & 由主题指定 & 渐变起点，或纯色填充的颜色。\\
\key{right color} & 由主题指定 & 渐变终点。\\
\key{track color} & 由主题指定 & 空轨道颜色。\\
\key{text color} & 由主题指定 & 标签文字颜色；auto 自动选黑白。\\
\key{font} & 沿用周围字体 & 标签字体声明。\\
\key{style} & 未指定 & 应用已定义的命名样式。\\
\key{missing text} & \key{--} & 空参数或 NA 的显示标签。\\
\key{number format} & \key{fixed} & fixed、grouped 或 scientific。\\
\key{band} & 空 & 范围内严格递增的两个参考端点。\\
\key{band color} & \key{gold} & 区间背景颜色。\\
\key{band opacity} & \key{0.3} & 区间透明度，范围为 0 到 1。\\
\key{stack colors} & 蓝、青绿、金黄 & 堆叠条分段颜色列表，可循环使用。\\
\bottomrule
\end{tabularx}
\endgroup


\subsection{数据、标签与图例参数}
\begin{tabularx}{\linewidth}{@{}p{36mm}p{23mm}X@{}}
\toprule
\textbf{参数} & \textbf{默认值} & \textbf{含义}\\
\midrule
\key{stack names} & 空 & 与各段一一对应的名称列表。\\
\key{legend} & \key{false} & 自动在堆叠条下方生成图例。\\
\key{segment labels} & \key{none} & none、value、percent、name、name value、name percent。\\
\key{segment font} & \key{\scriptsize} & 段标签字体。\\
\key{segment text color} & \key{auto} & 自动黑白或显式颜色。\\
\key{label overflow} & \key{auto} & 自动外移或 allow 允许越界。\\
\key{better} & \key{higher} & higher 越大越好，lower 越小越好；只影响条件配色。\\
\key{pattern} & \key{none} & 普通条的纹理名称。\\
\key{pattern color} & \key{black} & 纹理线条或点阵的颜色。\\
\key{stack patterns} & \key{none} & 堆叠段和图例的纹理列表，循环使用。\\
\key{print} & 未应用 & 黑白打印预设。\\
\key{print labels} & \key{false} & 内部标签白底；print 会开启此项。\\
\bottomrule
\end{tabularx}
\key{\gradbarslegend[选项]{名称列表}} 生成可共享的独立图例；
\key{\gradbarscolumn{G}{选项}} 定义统一样式的数值列。

CSV 命令包括 \key{\gradbarsloadcsv}、\key{\gradbarscsvtable}、\key{\gradbarscsvstyle}，\par
以及 \key{\gradbarcsv} 和 \key{\gradbarscsvcell}。
完整参数与范围计算规则见“CSV 数据导入”一节。

\subsection{图形与排版参数}
\begin{tabularx}{\linewidth}{@{}p{32mm}p{22mm}X@{}}
\toprule
\textbf{图形参数} & \textbf{默认值} & \textbf{说明}\\
\midrule
\key{shape} & \key{bar} & bar 普通条；lollipop 棒棒糖。哑铃与趋势使用专用命令。\\
\key{compare color} & \key{black!20} & 双层条的基准层颜色。\\
\key{compare ratio} & \key{.45} & 当前层高度比，严格在 0 与 1 之间。\\
\key{range point} & 空 & 浮动区间内的点，仅用于 gradrange。\\
\key{marker size} & \key{2pt} & 圆点半径，必须为正，受轨道高度限制。\\
\key{stem width} & \key{.6pt} & 棒棒糖及区间端点线宽，必须为正。\\
\key{outline} & \key{false} & 是否只绘制轮廓。\\
\key{outline width} & \key{.4pt} & 轮廓线宽，必须为正。\\
\key{row padding} & \key{0pt} & 图形上下各增加的留白，必须非负。\\
\key{preset} & 未应用 & paper、report、presentation、outline。\\
\bottomrule
\end{tabularx}


\subsection{图形与评价接口}
\begin{tabularx}{\linewidth}{@{}>{\raggedright\arraybackslash}p{43mm}X@{}}
\toprule
接口或选项 & 说明\\\midrule
\key{\graddumbbell[o]{a}{b}} & 基准 a、当前 b；空心与实心圆。\\
\key{compare label} & values（默认）、delta 或 percent；也适用于双层条。\\
\key{stack totals} & net（默认净值）或 separate（正小计 / 负小计）。\\
\key{\gradspark[o]{list}} & 等间距序列；空字段和 NA 为缺失。\\
\key{spark range} & auto（默认行内范围）或 fixed（公共 min/max）。\\
\key{spark points} & extrema（默认）、all、last、none。\\
\key{spark high color} & 默认 gradbarsTeal。\\
\key{spark low color} & 默认 gradbarsOrange。\\
\key{better} & higher、lower、target、interval。\\
\key{quality target} & target 模式的数值目标，默认空，必须提供。\\
\key{quality range} & interval 模式的闭区间，默认空，必须提供。\\
\key{palette} & categorical、sequential、diverging、mono，默认不应用。\\
\key{\gradbarscategory{n}{c}{p}} & 固定名称 n 的颜色 c 和纹理 p，定义遵循作用域。\\
\key{references} & 多条参考记录：数值/颜色/线型/名称；默认空。\\
\key{\gradbullet[o]{v}} & 分级背景与居中的实际值条，需设置 bullet bands。\\
\key{\gradintervals[o]{list}} & 多组区间：名称/下界/点估计/上界/颜色。\\
\bottomrule
\end{tabularx}

\clearpage
\subsection{表格与标注接口}
\begin{tabularx}{\linewidth}{@{}lX@{}}
\toprule
接口 & 含义\\\midrule
\key{gradtable} & 表格环境：外观、列宽与分页设置。\\
\key{\gradcolumn[o]{id}{title}} & 声明列：type 选择类型，options 接受图形与数字选项。\\
\key{\gradrow{list}} & 按列顺序填写逗号分隔的数据。\\
\key{\gradheader{spans}} & 添加一层表头，用 gradspan 声明连续的合并列。\\
\key{\gradspan{n}{title}} & 在 gradheader 中合并 n 列并设置标题。\\
\key{\gradgroup{title}} & 插入跨列分组标题，与首条数据保持在同页。\\
\key{\gradtablestyle{n}{keys}} & 保存表格选项，通过 style=n 调用。\\
\key{\gradtableschema{n}{body}} & 保存列定义，通过 graduseschema 调用。\\
\key{mark=none|best|both} & 数字、普通条与子弹图列的最佳及次佳值标记。\\
\key{width / weight} & 表格总宽度，以及可伸缩列分配剩余空间的权重。\\
\key{long=true} & 单栏正文中的跨页模式，重复完整表头。\\
\key{overflow mark} & none（单图默认）、arrow、slash。\\
\key{marker shape} & circle（默认）、square、diamond、triangle。\\
\key{marker fill} & solid（默认）或 hollow。\\
\key{compare annotation} & none（默认）、arrow、bracket。\\
\key{compare label=percent} & 相对于基准绝对值的百分比变化。\\
\bottomrule
\end{tabularx}

\subsection{已知边界}
CSV 仅支持逗号分隔的单行字段；数值列可自动计算共享范围。
本版不提供 CSV 字段内换行或 CSV 表格自动分页。趋势线使用等间距观测，不解析日期。
不要把绘图命令直接放入 siunitx 的数值型 S 列或 PDF 书签。
宏包已通过 pdfLaTeX、XeLaTeX、LuaLaTeX 的功能检查（包括 CSV）。本中文手册采用 XeLaTeX，英文手册可用 pdfLaTeX 编译。图形尚未生成专用的无障碍 PDF 语义标签。

\end{document}
